% Chapter 7. All quoted numbers come from chapters/ch07-numbers.tex, written by code/ch07/run_all.py.
% Generated by code/ch07/run_all.py. Do not edit by hand.
\newcommand{\alHiTecA}{5.25}
\newcommand{\alHiTecCapm}{0.28}
\newcommand{\alHiTecT}{4.30}
\newcommand{\alHlthA}{2.52}
\newcommand{\alHlthIR}{0.23}
\newcommand{\alHlthResid}{10.9}
\newcommand{\alHlthSe}{1.38}
\newcommand{\alHlthT}{1.89}
\newcommand{\alNW}{6}
\newcommand{\alNsigCapm}{0}
\newcommand{\alNsigCapmWord}{zero}
\newcommand{\alNsigSix}{3}
\newcommand{\alNsigSixWord}{three}
\newcommand{\alOtherA}{\ensuremath{-}2.23}
\newcommand{\alOtherT}{\ensuremath{-}3.02}
\newcommand{\alRsqMaxCapm}{0.85}
\newcommand{\alRsqMaxSix}{0.91}
\newcommand{\alRsqMinCapm}{0.33}
\newcommand{\alRsqMinSix}{0.43}
\newcommand{\alTableRows}{NoDur & 2.20 & 1.78 & 0.76 & 0.65 & \ensuremath{-}1.29 & \ensuremath{-}1.33 & 0.75 \\
Durbl & \ensuremath{-}1.02 & \ensuremath{-}0.50 & 1.25 & 0.56 & 0.42 & 0.21 & 0.60 \\
Manuf & 0.00 & 0.00 & 1.04 & 0.85 & \ensuremath{-}2.13 & \ensuremath{-}2.80 & 0.88 \\
Enrgy & 2.23 & 0.98 & 0.84 & 0.38 & \ensuremath{-}2.04 & \ensuremath{-}1.02 & 0.46 \\
HiTec & 0.28 & 0.19 & 1.24 & 0.75 & 5.25 & 4.30 & 0.84 \\
Telcm & \ensuremath{-}0.27 & \ensuremath{-}0.18 & 0.77 & 0.53 & 0.28 & 0.19 & 0.58 \\
Shops & 1.23 & 1.02 & 1.00 & 0.75 & 0.34 & 0.31 & 0.79 \\
Hlth & 2.52 & 1.89 & 0.81 & 0.57 & 1.52 & 1.13 & 0.62 \\
Utils & 1.92 & 1.34 & 0.52 & 0.33 & \ensuremath{-}1.19 & \ensuremath{-}0.89 & 0.43 \\
Other & \ensuremath{-}0.70 & \ensuremath{-}0.69 & 1.09 & 0.85 & \ensuremath{-}2.23 & \ensuremath{-}3.02 & 0.91 \\}
\newcommand{\atActive}{3.30}
\newcommand{\atAlloc}{1.39}
\newcommand{\atAllocHiTec}{\ensuremath{-}0.24}
\newcommand{\atB}{12.64}
\newcommand{\atBSum}{99.8}
\newcommand{\atBestSel}{HiTec}
\newcommand{\atBestSelV}{0.99}
\newcommand{\atEffSum}{3.31}
\newcommand{\atHiTecW}{42.2}
\newcommand{\atInter}{\ensuremath{-}1.04}
\newcommand{\atKmax}{1.194}
\newcommand{\atKmin}{1.041}
\newcommand{\atLast}{August 2026}
\newcommand{\atMonths}{8}
\newcommand{\atMonthsWord}{eight}
\newcommand{\atP}{15.94}
\newcommand{\atSel}{2.96}
\newcommand{\atSelHiTec}{0.99}
\newcommand{\atSumMonthly}{2.97}
\newcommand{\atTableRows}{NoDur & 2.4 & 2.4 & 0.03 & \ensuremath{-}0.24 & 0.02 \\
Durbl & 2.8 & 2.8 & 0.07 & 0.65 & \ensuremath{-}0.09 \\
Manuf & 6.9 & 6.9 & 0.23 & 0.68 & \ensuremath{-}0.17 \\
Enrgy & 2.4 & 5.4 & 1.11 & \ensuremath{-}0.06 & \ensuremath{-}0.17 \\
HiTec & 42.2 & 32.2 & \ensuremath{-}0.24 & 0.99 & \ensuremath{-}0.32 \\
Telcm & 1.7 & 1.7 & 0.02 & 0.33 & \ensuremath{-}0.07 \\
Shops & 10.1 & 10.1 & 0.08 & 0.13 & \ensuremath{-}0.03 \\
Hlth & 8.1 & 12.1 & 0.16 & \ensuremath{-}0.20 & \ensuremath{-}0.21 \\
Utils & 2.5 & 5.5 & \ensuremath{-}0.14 & 0.03 & 0.04 \\
Other & 20.7 & 20.7 & 0.07 & 0.65 & \ensuremath{-}0.04 \\
\midrule
Total & 100.0 & 100.0 & 1.39 & 2.96 & \ensuremath{-}1.04 \\}
\newcommand{\atWSum}{99.8}
\newcommand{\atWorstSel}{NoDur}
\newcommand{\atWorstSelV}{\ensuremath{-}0.24}
\newcommand{\blBlSum}{1.042}
\newcommand{\blDelta}{3.03}
\newcommand{\blLamOne}{0.271}
\newcommand{\blLamTwo}{0.042}
\newcommand{\blMuEnrgy}{4.12}
\newcommand{\blMuHiTec}{7.28}
\newcommand{\blMuHlth}{5.30}
\newcommand{\blMvGross}{12.7}
\newcommand{\blMvMax}{254}
\newcommand{\blMvMaxName}{Telcm}
\newcommand{\blMvSum}{1.36}
\newcommand{\blPiEnrgy}{2.68}
\newcommand{\blPiHiTec}{10.33}
\newcommand{\blPiHlth}{4.28}
\newcommand{\blPiMax}{14.37}
\newcommand{\blPiMaxName}{Durbl}
\newcommand{\blPiMin}{2.68}
\newcommand{\blPiMinName}{Enrgy}
\newcommand{\blTau}{0.0167}
\newcommand{\blTauFrac}{1/60}
\newcommand{\blViewOneAfter}{\ensuremath{-}1.98}
\newcommand{\blViewOneBefore}{\ensuremath{-}6.05}
\newcommand{\blViewTwoAfter}{4.12}
\newcommand{\blViewTwoBefore}{2.68}
\newcommand{\blW}{60}
\newcommand{\blWblEnrgy}{7.2}
\newcommand{\blWblHiTec}{15.7}
\newcommand{\blWblHlth}{34.9}
\newcommand{\blWmEnrgy}{3.0}
\newcommand{\blWmHiTec}{42.8}
\newcommand{\blWmHlth}{7.8}
\newcommand{\blWmvEnrgy}{195}
\newcommand{\blWmvHiTec}{253}
\newcommand{\blWmvHlth}{66}
\newcommand{\cpTanBetaMax}{1.03}
\newcommand{\cpTanBetaMin}{0.59}
\newcommand{\csBest}{minimum variance, shrunk}
\newcommand{\csCapA}{0.14}
\newcommand{\csCapAt}{2.41}
\newcommand{\csCapB}{0.99}
\newcommand{\csCapMean}{7.23}
\newcommand{\csCapP}{0.14}
\newcommand{\csCapSR}{0.46}
\newcommand{\csCapTurn}{0}
\newcommand{\csCapVol}{15.7}
\newcommand{\csCapWorst}{23.1}
\newcommand{\csCapZ}{\ensuremath{-}1.49}
\newcommand{\csEnd}{August 2026}
\newcommand{\csEwA}{0.84}
\newcommand{\csEwAt}{1.87}
\newcommand{\csEwB}{0.93}
\newcommand{\csEwMean}{7.46}
\newcommand{\csEwSR}{0.50}
\newcommand{\csEwTurn}{2}
\newcommand{\csEwVol}{15.0}
\newcommand{\csEwWorst}{21.7}
\newcommand{\csFallback}{4}
\newcommand{\csGmvA}{2.97}
\newcommand{\csGmvAt}{2.28}
\newcommand{\csGmvB}{0.49}
\newcommand{\csGmvLev}{13.2}
\newcommand{\csGmvMean}{6.45}
\newcommand{\csGmvP}{0.99}
\newcommand{\csGmvRes}{64}
\newcommand{\csGmvSR}{0.50}
\newcommand{\csGmvTurn}{32}
\newcommand{\csGmvVol}{12.9}
\newcommand{\csGmvWorst}{13.0}
\newcommand{\csGmvZ}{0.02}
\newcommand{\csLogmvA}{1.98}
\newcommand{\csLogmvAt}{2.04}
\newcommand{\csLogmvB}{0.65}
\newcommand{\csLogmvLev}{10.2}
\newcommand{\csLogmvMean}{6.59}
\newcommand{\csLogmvP}{0.72}
\newcommand{\csLogmvRes}{34}
\newcommand{\csLogmvSR}{0.52}
\newcommand{\csLogmvTurn}{10}
\newcommand{\csLogmvVol}{12.6}
\newcommand{\csLogmvWorst}{12.9}
\newcommand{\csLogmvZ}{0.36}
\newcommand{\csLotanA}{1.27}
\newcommand{\csLotanAt}{1.18}
\newcommand{\csLotanB}{0.87}
\newcommand{\csLotanMean}{7.45}
\newcommand{\csLotanP}{0.63}
\newcommand{\csLotanSR}{0.46}
\newcommand{\csLotanTurn}{26}
\newcommand{\csLotanVol}{16.2}
\newcommand{\csLotanWorst}{19.1}
\newcommand{\csLotanZ}{\ensuremath{-}0.48}
\newcommand{\csLwA}{2.60}
\newcommand{\csLwAt}{2.53}
\newcommand{\csLwB}{0.57}
\newcommand{\csLwLev}{11.7}
\newcommand{\csLwMean}{6.69}
\newcommand{\csLwP}{0.61}
\newcommand{\csLwRes}{46}
\newcommand{\csLwSR}{0.54}
\newcommand{\csLwTurn}{17}
\newcommand{\csLwVol}{12.3}
\newcommand{\csLwWorst}{12.5}
\newcommand{\csLwZ}{0.51}
\newcommand{\csMktMean}{7.13}
\newcommand{\csShrinkMed}{0.09}
\newcommand{\csStart}{July 1968}
\newcommand{\csTT}{698}
\newcommand{\csTableRows}{Equal weights & 7.46 & 15.0 & 0.50 &  & 2 & 0.84 & 0.93 \\
Market weights & 7.23 & 15.7 & 0.46 & \ensuremath{-}1.49 & 0 & 0.14 & 0.99 \\
Minimum variance & 6.45 & 12.9 & 0.50 & 0.02 & 32 & 2.97 & 0.49 \\
Minimum variance, shrunk & 6.69 & 12.3 & 0.54 & 0.51 & 17 & 2.60 & 0.57 \\
Minimum variance, long only & 6.59 & 12.6 & 0.52 & 0.36 & 10 & 1.98 & 0.65 \\
Tangency, scaled & 2.53 & 20.8 & 0.12 & \ensuremath{-}2.24 & 118 & 0.35 & 0.31 \\
Tangency, long only & 7.45 & 16.2 & 0.46 & \ensuremath{-}0.48 & 26 & 1.27 & 0.87 \\}
\newcommand{\csTanA}{0.35}
\newcommand{\csTanAt}{0.14}
\newcommand{\csTanB}{0.31}
\newcommand{\csTanMean}{2.53}
\newcommand{\csTanP}{0.03}
\newcommand{\csTanSR}{0.12}
\newcommand{\csTanTurn}{118}
\newcommand{\csTanVol}{20.8}
\newcommand{\csTanWorst}{35.6}
\newcommand{\csTanZ}{\ensuremath{-}2.24}
\newcommand{\csTenEwSR}{0.55}
\newcommand{\csTenLotanSR}{0.45}
\newcommand{\csTenLwSR}{0.60}
\newcommand{\csTenLwZ}{0.49}
\newcommand{\csTenTanSR}{0.19}
\newcommand{\csTenTanZ}{\ensuremath{-}2.21}
\newcommand{\csW}{60}
\newcommand{\cvInMax}{0.79}
\newcommand{\cvM}{2,000}
\newcommand{\cvOutMax}{0.53}
\newcommand{\cvShareMax}{82}
\newcommand{\cvTmax}{600}
\newcommand{\dtSRm}{0.47}
\newcommand{\dtYrsHalf}{16}
\newcommand{\dtYrsOne}{4}
\newcommand{\dtYrsQuarter}{65}
\newcommand{\dvLimit}{14.2}
\newcommand{\dvOne}{18.7}
\newcommand{\dvRhoBar}{0.60}
\newcommand{\dvShare}{38}
\newcommand{\dvShareInf}{43}
\newcommand{\dvTen}{14.7}
\newcommand{\esInOneTwenty}{1.25}
\newcommand{\esInSixty}{1.73}
\newcommand{\esInTheoryOneTwenty}{1.25}
\newcommand{\esInTheorySixty}{1.72}
\newcommand{\esInTheoryTwoForty}{0.98}
\newcommand{\esInTwoForty}{0.98}
\newcommand{\esM}{4,000}
\newcommand{\esNegOneTwenty}{2}
\newcommand{\esNegSixty}{9}
\newcommand{\esNegTwoForty}{0}
\newcommand{\esOutOneTwenty}{0.33}
\newcommand{\esOutSixty}{0.24}
\newcommand{\esOutTwoForty}{0.43}
\newcommand{\esTheta}{0.64}
\newcommand{\faAlpha}{5.25}
\newcommand{\faBCma}{\ensuremath{-}0.41}
\newcommand{\faBHml}{\ensuremath{-}0.38}
\newcommand{\faBMkt}{1.06}
\newcommand{\faBMom}{\ensuremath{-}0.06}
\newcommand{\faBRmw}{\ensuremath{-}0.31}
\newcommand{\faBSmb}{0.07}
\newcommand{\faCapmPred}{8.92}
\newcommand{\faCma}{\ensuremath{-}1.18}
\newcommand{\faFourNeg}{3.87}
\newcommand{\faFourNegSigned}{\ensuremath{-}3.87}
\newcommand{\faFourRoundedNeg}{\ensuremath{-}3.86}
\newcommand{\faFourRoundedSum}{3.86}
\newcommand{\faGap}{4.96}
\newcommand{\faHml}{\ensuremath{-}1.34}
\newcommand{\faMean}{9.20}
\newcommand{\faMkt}{7.67}
\newcommand{\faMktDrop}{1.25}
\newcommand{\faMom}{\ensuremath{-}0.41}
\newcommand{\faRmw}{\ensuremath{-}0.93}
\newcommand{\faRoundedSum}{9.22}
\newcommand{\faSixPred}{3.96}
\newcommand{\faSmb}{0.16}
\newcommand{\fmCorrErr}{0.12}
\newcommand{\fmGapMax}{\ensuremath{-}0.12}
\newcommand{\fmGapPair}{Enrgy and Shops}
\newcommand{\fmGapPairs}{45}
\newcommand{\fmGapPos}{17}
\newcommand{\fmGmvFactor}{11.48}
\newcommand{\fmGmvFactorTrue}{12.41}
\newcommand{\fmGmvSample}{12.05}
\newcommand{\fqFirst}{July 1963}
\newcommand{\fqLast}{August 2026}
\newcommand{\fqT}{758}
\newcommand{\frCapSR}{0.48}
\newcommand{\frGmvMu}{6.72}
\newcommand{\frGmvVol}{12.1}
\newcommand{\frGross}{2.95}
\newcommand{\frLoGmvVol}{12.4}
\newcommand{\frLoN}{6}
\newcommand{\frLoNword}{six}
\newcommand{\frLoSR}{0.58}
\newcommand{\frMktMu}{7.22}
\newcommand{\frMktSR}{0.47}
\newcommand{\frMktVol}{15.4}
\newcommand{\frShortMax}{55}
\newcommand{\frShortName}{Other}
\newcommand{\frTableRows}{NoDur & 27.8 & 57.1 & 16.3 & 28.7 & 2.5 \\
Durbl & 0.8 & 4.7 & 0.0 & 0.0 & 2.3 \\
Manuf & \ensuremath{-}1.8 & \ensuremath{-}29.2 & 0.0 & 0.0 & 7.8 \\
Enrgy & 7.1 & 27.0 & 1.8 & 14.2 & 3.0 \\
HiTec & 4.5 & 24.6 & 0.0 & 8.8 & 42.8 \\
Telcm & 22.6 & \ensuremath{-}13.3 & 19.6 & 0.0 & 1.4 \\
Shops & 8.3 & 28.5 & 0.0 & 5.8 & 10.1 \\
Hlth & 19.4 & 33.4 & 16.4 & 27.7 & 7.8 \\
Utils & 45.5 & 22.3 & 45.9 & 14.9 & 2.5 \\
Other & \ensuremath{-}34.3 & \ensuremath{-}55.1 & 0.0 & 0.0 & 19.7 \\
\midrule
Mean, \% a year & 6.72 & 8.90 & 6.42 & 7.84 & 8.29 \\
Volatility, \% a year & 12.1 & 13.9 & 12.4 & 13.5 & 17.2 \\
Sharpe ratio & 0.56 & 0.64 & 0.52 & 0.58 & 0.48 \\}
\newcommand{\frTanMu}{8.90}
\newcommand{\frTanSR}{0.64}
\newcommand{\frTanVol}{13.9}
\newcommand{\fwDoubleGross}{11.3}
\newcommand{\fwDoubleMu}{17.8}
\newcommand{\fwDoubleShort}{160}
\newcommand{\fwDoubleShortName}{Telcm}
\newcommand{\grsCapm}{1.39}
\newcommand{\grsFive}{4.02}
\newcommand{\grsPCapm}{0.18}
\newcommand{\grsPFive}{\ensuremath{2.2\times 10^{-5}}}
\newcommand{\grsPSix}{\ensuremath{1.2\times 10^{-4}}}
\newcommand{\grsSRall}{1.40}
\newcommand{\grsSRcapm}{0.47}
\newcommand{\grsSRcapmAll}{0.67}
\newcommand{\grsSRfac}{1.15}
\newcommand{\grsSix}{3.56}
\newcommand{\mgBh}{0}
\newcommand{\mgBhFalse}{0}
\newcommand{\mgBhTrue}{0}
\newcommand{\mgLongBh}{16}
\newcommand{\mgLongBhFalse}{2}
\newcommand{\mgLongNaive}{97}
\newcommand{\mgLongNaiveFalse}{42}
\newcommand{\mgLongPower}{60}
\newcommand{\mgN}{1,000}
\newcommand{\mgNaive}{85}
\newcommand{\mgNaiveFalse}{53}
\newcommand{\mgNaivePos}{62}
\newcommand{\mgNaivePosFalse}{30}
\newcommand{\mgNaiveTrue}{32}
\newcommand{\mgPiZero}{0.95}
\newcommand{\mgPowerTheory}{35}
\newcommand{\mgSkill}{100}
\newcommand{\mgT}{120}
\newcommand{\pcMktCorr}{0.98}
\newcommand{\pcOne}{65}
\newcommand{\pcThree}{80}
\newcommand{\pcTwo}{9}
\newcommand{\pcTwoHml}{0.39}
\newcommand{\pcTwoHmlSigned}{\ensuremath{-}0.39}
\newcommand{\pcTwoMkt}{0.09}
\newcommand{\poAtmPayoff}{1.46}
\newcommand{\poAtmPrice}{1.67}
\newcommand{\poFirst}{July 1964}
\newcommand{\poMarkupAdj}{0.15}
\newcommand{\poMarkupCapm}{2.09}
\newcommand{\poMoney}{95}
\newcommand{\poT}{746}
\newcommand{\poTwentyA}{0.64}
\newcommand{\poTwentyAdj}{\ensuremath{-}0.75}
\newcommand{\poTwentyB}{0.18}
\newcommand{\poTwentyCapmR}{0.34}
\newcommand{\poTwentyFlat}{4.21}
\newcommand{\poTwentyFlatN}{58}
\newcommand{\poTwentyHmA}{10.72}
\newcommand{\poTwentyHmB}{\ensuremath{-}0.06}
\newcommand{\poTwentyHmG}{\ensuremath{-}0.48}
\newcommand{\poTwentyHmR}{0.59}
\newcommand{\poTwentyMean}{1.93}
\newcommand{\poTwentyPrem}{5.33}
\newcommand{\poTwentySR}{0.40}
\newcommand{\poTwentyT}{0.91}
\newcommand{\poTwentyTmG}{\ensuremath{-}2.30}
\newcommand{\poTwentyTmT}{\ensuremath{-}11.05}
\newcommand{\poTwentyVol}{4.8}
\newcommand{\poTwentyWorst}{18.2}
\newcommand{\poTwentyWorstWhen}{October 1987}
\newcommand{\poZeroA}{\ensuremath{-}1.44}
\newcommand{\poZeroAdj}{\ensuremath{-}0.90}
\newcommand{\poZeroB}{0.18}
\newcommand{\poZeroCapmR}{0.35}
\newcommand{\poZeroHmA}{9.07}
\newcommand{\poZeroHmB}{\ensuremath{-}0.07}
\newcommand{\poZeroHmG}{\ensuremath{-}0.50}
\newcommand{\poZeroHmR}{0.63}
\newcommand{\poZeroMean}{\ensuremath{-}0.18}
\newcommand{\poZeroPrem}{3.23}
\newcommand{\poZeroSR}{\ensuremath{-}0.04}
\newcommand{\poZeroT}{\ensuremath{-}2.10}
\newcommand{\poZeroTmG}{\ensuremath{-}2.36}
\newcommand{\poZeroTmT}{\ensuremath{-}11.69}
\newcommand{\poZeroVol}{4.7}
\newcommand{\poZeroWorst}{18.4}
\newcommand{\poZeroWorstWhen}{October 1987}
\newcommand{\rbDivActual}{0.66}
\newcommand{\rbDivApprox}{0.67}
\newcommand{\rbGrowAvg}{10.14}
\newcommand{\rbGrowHold}{10.41}
\newcommand{\rbGrowReb}{10.79}
\newcommand{\rbHoldTop}{18}
\newcommand{\rbHoldTopName}{Hlth}
\newcommand{\rbLastEnrgy}{0.37}
\newcommand{\rbLastHiTec}{1.26}
\newcommand{\rbLastUtils}{0.55}
\newcommand{\rbMaxEnrgy}{1.73}
\newcommand{\rbMaxHiTec}{2.02}
\newcommand{\rbMaxUtils}{0.81}
\newcommand{\rbMaxWhenEnrgy}{February 2021}
\newcommand{\rbMaxWhenHiTec}{October 2004}
\newcommand{\rbMaxWhenUtils}{August 1974}
\newcommand{\rbMinEnrgy}{0.36}
\newcommand{\rbMinHiTec}{1.00}
\newcommand{\rbMinUtils}{0.01}
\newcommand{\rbMinWhenEnrgy}{July 2026}
\newcommand{\rbMinWhenHiTec}{July 2021}
\newcommand{\rbMinWhenUtils}{February 2001}
\newcommand{\rbOutside}{45}
\newcommand{\rbSeEnrgy}{0.12}
\newcommand{\rbSeHiTec}{0.08}
\newcommand{\rbSeUtils}{0.09}
\newcommand{\rbW}{60}
\newcommand{\rbWealthHold}{718}
\newcommand{\rbWealthReb}{915}
\newcommand{\rpEqMktShare}{25}
\newcommand{\rpErcCma}{20.9}
\newcommand{\rpErcHml}{12.7}
\newcommand{\rpErcMax}{23.9}
\newcommand{\rpErcMaxName}{RMW}
\newcommand{\rpErcMin}{12.7}
\newcommand{\rpErcMinName}{HML}
\newcommand{\rpErcMkt}{13.7}
\newcommand{\rpErcMom}{12.8}
\newcommand{\rpLev}{4.2}
\newcommand{\rpSReq}{1.08}
\newcommand{\rpSRerc}{1.09}
\newcommand{\rpSRmkt}{0.47}
\newcommand{\rpSRmktSame}{0.45}
\newcommand{\rpSRrolling}{1.10}
\newcommand{\rpVolErc}{3.7}
\newcommand{\smBetaMax}{1.25}
\newcommand{\smBetaMaxName}{Durbl}
\newcommand{\smBetaMin}{0.52}
\newcommand{\smBetaMinName}{Utils}
\newcommand{\smIcpt}{4.51}
\newcommand{\smRsq}{0.38}
\newcommand{\smSlope}{3.28}
\newcommand{\smSlopeT}{2.21}
\newcommand{\smoAcfEight}{0.81}
\newcommand{\smoAcfHalf}{0.54}
\newcommand{\smoAlpha}{\ensuremath{-}0.70}
\newcommand{\smoAlphaEight}{5.41}
\newcommand{\smoAlphaHalf}{3.13}
\newcommand{\smoAlphaTEight}{8.04}
\newcommand{\smoAlphaTHalf}{3.66}
\newcommand{\smoBeta}{1.09}
\newcommand{\smoBetaEight}{0.22}
\newcommand{\smoBetaHalf}{0.55}
\newcommand{\smoDdEight}{51}
\newcommand{\smoDdFirst}{October 2007}
\newcommand{\smoDdHalf}{60}
\newcommand{\smoDdLast}{December 2009}
\newcommand{\smoDdTrue}{67}
\newcommand{\smoDimsonAlphaEight}{\ensuremath{-}0.12}
\newcommand{\smoDimsonAlphaHalf}{\ensuremath{-}0.63}
\newcommand{\smoDimsonEight}{1.01}
\newcommand{\smoDimsonHalf}{1.08}
\newcommand{\smoDimsonSixEight}{0.88}
\newcommand{\smoDimsonSixHalf}{1.13}
\newcommand{\smoMeanKept}{99.2}
\newcommand{\smoMktAcf}{0.04}
\newcommand{\smoOwnAcf}{0.08}
\newcommand{\smoQ}{12}
\newcommand{\smoSR}{0.39}
\newcommand{\smoSREight}{1.11}
\newcommand{\smoSRFormula}{0.68}
\newcommand{\smoSRHalf}{0.65}
\newcommand{\smoSpurEight}{0.58}
\newcommand{\smoSpurHalf}{0.06}
\newcommand{\smoVol}{18.3}
\newcommand{\smoVolEight}{6.3}
\newcommand{\smoVolExact}{10.9}
\newcommand{\smoVolFormula}{10.5}
\newcommand{\smoVolHalf}{10.9}
\newcommand{\xrEightMuA}{4.00}
\newcommand{\xrEightPiA}{5.0}
\newcommand{\xrEightPiB}{11.25}
\newcommand{\xrEightWA}{40.0}
\newcommand{\xrEightWB}{50.0}
\newcommand{\xrFiveA}{3.6}
\newcommand{\xrFiveB}{0.60}
\newcommand{\xrFiveD}{1.05}
\newcommand{\xrFour}{25.4}
\newcommand{\xrNineVol}{6.20}
\newcommand{\xrNineW}{20}
\newcommand{\xrOneVol}{9.79}
\newcommand{\xrOneW}{10.53}
\newcommand{\xrSevenB}{2.01}
\newcommand{\xrSevenKone}{1.0200}
\newcommand{\xrSevenKtwo}{1.0150}
\newcommand{\xrSevenOne}{1.0200}
\newcommand{\xrSevenP}{5.06}
\newcommand{\xrSevenSum}{3.00}
\newcommand{\xrSevenTwo}{2.0300}
\newcommand{\xrSixActive}{1.90}
\newcommand{\xrSixAlloc}{0.80}
\newcommand{\xrSixInt}{0.30}
\newcommand{\xrSixRb}{6.80}
\newcommand{\xrSixSel}{0.80}
\newcommand{\xrThreeA}{0.722}
\newcommand{\xrThreeB}{1.083}
\newcommand{\xrThreeC}{1.444}
\newcommand{\xrTwoA}{46.2}
\newcommand{\xrTwoB}{30.8}
\newcommand{\xrTwoC}{23.1}
\newcommand{\xrTwoMu}{5.54}
\newcommand{\xrTwoSR}{0.693}
\newcommand{\xrTwoVol}{7.99}

\begin{refsection}
\setcounter{chapter}{6}
\chapter{Portfolio Theory, Factor Models, Alpha and Beta}
\label{ch:portfolio}

An allocator who adds a hedge fund, a private credit sleeve or a timber fund to a pension plan is
making two claims at once. The first is that the new asset improves the whole portfolio, which is a
statement about means, variances and covariances. The second is that the manager earns something
the plan could not get more cheaply by holding index funds, which is a statement about alpha and
beta. Both claims rest on the same algebra, and both are much harder to verify from data than the
algebra suggests. This chapter builds the algebra, measures how far the data can be trusted, and
ends with a test on more than half a century of returns: would an investor who estimated the
mean-variance frontier from the previous five years have done better than one who simply held equal
weights?

The first half is portfolio choice. Diversification by equal weights has a floor set by the average covariance. The
frontier of portfolios with the least variance for a given mean has a closed form, and adding a
riskless asset picks out one portfolio on it, the tangency portfolio, which every mean-variance
investor with the same means and covariances holds in some proportion when all can borrow and lend
at the riskless rate. The CAPM is the statement that the market is that portfolio, and
it turns out to be equivalent to a linear relation between expected returns and betas. Factor models
generalise the relation and give a covariance matrix that can be estimated. The second half is
measurement. Alpha is the intercept of a regression, and its standard error, the number of years it
takes to detect, the joint test across many assets and the problem of screening many managers at
once all follow from least squares. Returns of illiquid assets are smoothed, which hides beta and
manufactures alpha, and option-like strategies produce alpha in a linear regression that is only a
premium for selling insurance. Attribution splits a manager's active return into decisions, the
Black-Litterman model blends equilibrium returns with views, and risk budgets allocate volatility
instead of capital.

\begin{keybox}[How to read this chapter]
The numbers in the text are computed by ten scripts in \texttt{code/ch07/} (\texttt{frontier},
\texttt{capm\_factors}, \texttt{alpha\_tests}, \texttt{smoothing\_beta}, \texttt{option\_alpha},
\texttt{attribution}, \texttt{black\_litterman}, \texttt{allocation}, \texttt{case\_oos} and
\texttt{exercises}); the tests in \texttt{code/ch07/test\_*.py} check them. An independent R script
recomputes the frontier, searches for the long-only tangency portfolio on its own, reruns every
regression with \texttt{lm}, the Newey-West errors and the joint test, applies the false discovery
rule with \texttt{p.adjust}, smooths and unsmooths the returns, reprices the put-writing fund, rebuilds
the attribution, the Black-Litterman posterior and the risk budgets, and recomputes the case study.
GNU Octave solves the long-only problems with its quadratic programming solver and the risk budget
with \texttt{fsolve}. The short listings printed in the text are in \texttt{code/ch07/snippets/};
\texttt{code/ch07/test\_assets.py} runs each one and checks the numbers it prints.

The data are monthly returns from the Kenneth R. French Data Library \parencite{frenchlib}: the market
excess return, the size, value, profitability, investment and momentum factors, the one month
Treasury bill rate, and the value- and equal-weighted returns and market values of ten industry
portfolios of US stocks. The sample runs from \fqFirst{} to \fqLast, \fqT{} months; \fqLast{} is
the last month in the library when the files were retrieved. Returns are simple monthly returns, and
excess returns are over the bill. Means and volatilities are quoted at annual rates, twelve times the
monthly mean and $\sqrt{12}$ times the monthly standard deviation; Sharpe and information ratios are
annualised the same way.

Random draws come from NumPy's default generator with base seed \texttt{707} and a stream number
fixed in each script. The seed is an arbitrary label that makes the runs repeatable; it carries no
date or other meaning. Simulated samples are drawn from normal distributions with the moments of the
data, so they show what estimation error does when the model is exactly right.

Notation. There are $N$ risky assets with return vector $R$, mean $\mu$ and covariance matrix
$\Sigma$, which is always assumed positive definite. $x=R-r_f\mathbf 1$ is the vector of excess
returns over the riskless rate $r_f$, with mean $\mu_x$, and $\mathbf 1$ is a vector of ones. A
portfolio is a weight vector $w$, fully invested when $\mathbf 1^\top w=1$. $\beta$ is a regression
slope and $\alpha$ an intercept; $\varepsilon$ is a regression residual with variance $\sigma_\varepsilon^2$.
$f$ is a vector of $K$ factors with covariance $\Omega$, and $B$ the $N\times K$ matrix of loadings.
$\theta$ is a Sharpe ratio and $\theta^2$ the largest squared Sharpe ratio available from a set of
assets. $T$ is a number of months. In \cref{sec7:smoothing} $\phi$ is a smoothing parameter and in
\cref{sec7:bl} $\delta$ is risk aversion and $\tau$ a scale for the uncertainty of the prior.
Some letters do double duty. The scalars $A,B,C,D$ of \cref{eq7:abcd} are unrelated to the loadings
$B$ and the residual covariance $D$ of \cref{def7:factor}; in \cref{prop7:carino} $B_t$ and $B$ are
benchmark returns; in \cref{sec7:bl} $\Omega$ is the covariance of the errors in the views; and in the
joint test of alphas the residual covariance matrix is written $\Sigma_\varepsilon$.
The letter $q$ is a number of lags in \cref{sec7:smoothing}, a false discovery level in the screening
of managers and the vector of views in \cref{sec7:bl}. The letter $m$ is a target mean on the frontier,
the mean of the stacked returns in the proof of \cref{prop7:grs}, the number of p-values in the
screening of managers and the market excess return in \cref{sec7:smoothing,sec7:options}. The letter
$\lambda$ is a multiplier in \cref{prop7:frontier,prop7:longonly} and the vector of factor prices in
\cref{prop7:apt}. The smoothing weights $\vartheta_k$ are unrelated to the Sharpe ratio $\theta$.
The letter $\pi$ is the vector of equilibrium returns in \cref{sec7:bl} and the price $\pi_t$ of an
option in \cref{sec7:options}; $M$ is the constraint matrix of \cref{prop7:longonly} and the covariance
of the views in the proof of \cref{prop7:bl}, and $G$ means different things in the proofs of
\cref{prop7:woodbury,prop7:bl}. The letter $e$ is the noise in the true return in \cref{prop7:stale}
and the error of the views in \cref{sec7:bl}, and $k_t$ and $k$ in \cref{prop7:carino} are linking
factors, not indices.
\end{keybox}

\FloatBarrier
% =====================================================================================
\section{How much risk does holding more assets remove?}
\label{sec7:diversify}

A portfolio $w$ of assets with mean vector $\mu$ and covariance matrix $\Sigma$ has mean return
$w^\top\mu$ and variance $w^\top\Sigma w=\sum_i\sum_j w_iw_j\Sigma_{ij}$. Of the $N^2$ terms in the
double sum only $N$ are variances, so as a portfolio spreads over more assets the covariances take
over. The equally weighted portfolio makes the point exactly.

\begin{proposition}[The floor under diversification]\label{prop7:diversify}
Let $\bar v=\frac1N\sum_i\Sigma_{ii}$ be the average variance and $\bar c=\frac{2}{N(N-1)}\sum_{i<j}\Sigma_{ij}$
the average covariance of $N$ assets. The portfolio with weight $1/n$ in each of $n$ assets chosen
from the $N$ has a variance whose average over all $\binom Nn$ choices is
\begin{equation}\label{eq7:diversify}
  \frac{\bar v}{n}+\Bigl(1-\frac1n\Bigr)\bar c .
\end{equation}
In particular the equally weighted portfolio of all $N$ assets has variance $\bar v/N+(1-1/N)\bar c$.
Since $\bar v\ge\bar c$ for any covariance matrix, \cref{eq7:diversify} exceeds $\bar c$ by $(\bar
v-\bar c)/n$; if $N$ and $n$ grow with $\bar v$ and $\bar c$ held fixed, which forces $\bar c\ge0$
because the variance at $n=N$ cannot be negative, it falls to $\bar c$ from above.
\end{proposition}

\begin{proof}
For a given subset $s$ of size $n$ the variance is $n^{-2}\bigl(\sum_{i\in s}\Sigma_{ii}+\sum_{i\ne j\in s}\Sigma_{ij}\bigr)$.
Averaging over subsets, each asset lies in a fraction $n/N$ of them and each ordered pair of distinct
assets in a fraction $n(n-1)/(N(N-1))$. The average of the first sum is therefore $n\bar v$ and of the
second $n(n-1)\bar c$, and dividing by $n^2$ gives \cref{eq7:diversify}.
\end{proof}

\begin{figure}[!tb]
  \centering
  \includegraphics[width=0.84\textwidth]{ch07/fig01_diversify}
  \caption{Volatility of equal-weighted portfolios of $n$ of the ten industries, \fqFirst{} to
  \fqLast. The line is the square root of the average variance over all subsets of size $n$,
  given by \cref{eq7:diversify}; by Jensen's inequality it lies above the average of the subsets'
  volatilities. The band runs from the least to the most volatile subset. The dashed line is the floor
  $\sqrt{\bar c}$, which bounds only the average: single subsets of up to eight industries fall below it.}
  \label{fig7:diversify}
\end{figure}

The ten industry portfolios of the data library are themselves diversified portfolios of hundreds of
stocks, so they are a demanding test. Over \fqFirst{} to \fqLast{} the square root of the average variance of a
single industry is \dvOne\% a year and the average correlation between two industries is \dvRhoBar. Equal
weights in all ten bring the volatility down to \dvTen\%, removing \dvShare\% of the average
variance, and the floor $\sqrt{\bar c}$ is \dvLimit\%. \Cref{fig7:diversify} shows the average over
subsets, which matches \cref{eq7:diversify} to rounding error, and the range across subsets, which is
wide for small $n$: two industries can diversify each other well or hardly at all. The lesson for an
allocator is that the second and third assets do most of the work, and that an asset whose
correlation with the rest is high adds little however many others there are.


\FloatBarrier
% =====================================================================================
\section{Which portfolios are efficient?}
\label{sec7:frontier}

\textcite{markowitz1952} asked which portfolios an investor who cares only about the mean and the
variance of the return would choose. For each target mean $m$ the answer is the fully invested
portfolio with the least variance, and the set of these portfolios is the minimum-variance frontier.
Write
\begin{equation}\label{eq7:abcd}
  A=\mathbf 1^\top\Sigma^{-1}\mathbf 1,\qquad B=\mathbf 1^\top\Sigma^{-1}\mu,\qquad C=\mu^\top\Sigma^{-1}\mu,
  \qquad D=AC-B^2 .
\end{equation}
When $\mu$ is not a multiple of $\mathbf 1$, $D>0$ by the Cauchy-Schwarz inequality in the inner
product $\langle a,b\rangle=a^\top\Sigma^{-1}b$.

\begin{proposition}[The frontier in closed form]\label{prop7:frontier}
The fully invested portfolio with mean $m$ and the least variance is
\begin{equation}\label{eq7:frontier}
  w(m)=\Sigma^{-1}\bigl(\lambda\mathbf 1+\gamma\mu\bigr),\qquad \lambda=\frac{C-Bm}{D},\qquad \gamma=\frac{Am-B}{D},
\end{equation}
and its variance is $\sigma^2(m)=(Am^2-2Bm+C)/D$, a parabola in $m$. The global minimum-variance
portfolio is $w_g=\Sigma^{-1}\mathbf 1/A$, with mean $B/A$ and variance $1/A$.
\end{proposition}

\begin{proof}
The problem is to minimise $\frac12w^\top\Sigma w$ subject to $\mathbf 1^\top w=1$ and $\mu^\top w=m$.
The objective is strictly convex and the constraints are linear, so the stationary point of the
Lagrangian $\frac12w^\top\Sigma w-\lambda(\mathbf 1^\top w-1)-\gamma(\mu^\top w-m)$ is the unique
minimum. Stationarity gives $\Sigma w=\lambda\mathbf 1+\gamma\mu$, and the two constraints become
$A\lambda+B\gamma=1$ and $B\lambda+C\gamma=m$, whose solution is the pair in \cref{eq7:frontier}.
The variance is $w^\top\Sigma w=w^\top(\lambda\mathbf 1+\gamma\mu)=\lambda+\gamma m$, which simplifies
to the parabola. Its minimum over $m$ is at $m=B/A$, where $\gamma=0$ and $\lambda=1/A$.
\end{proof}

\begin{corollary}[Two-fund separation]\label{cor7:twofund}
The weights $w(m)$ are affine in $m$, so every frontier portfolio is a combination of any two
distinct frontier portfolios, with weights that add to one.
\end{corollary}

Plotted with volatility on the horizontal axis, the parabola becomes the upper and lower branches of
a hyperbola. Portfolios on the upper branch are efficient: no other portfolio has a higher mean for
the same variance. \Cref{cor7:twofund} is the first appearance of an idea that runs through the
chapter: when everyone uses the same moments, the choice of portfolio collapses to a choice among a
few funds.

\Cref{fig7:weights} draws the weights of the ten industries along the frontier, with means and
covariances estimated from \fqFirst{} to \fqLast. Each weight is a straight line in
$m$, as the corollary says, so a target mean far from that of the minimum-variance portfolio is reached
only with large positions on both sides. At twice the mean of the tangency portfolio introduced below,
\fwDoubleMu\% a year, the frontier portfolio sells short \fwDoubleShort\% of its capital in
\fwDoubleShortName{} and holds gross positions of \fwDoubleGross{} times its capital.

\begin{figure}[!tb]
  \centering
  \includegraphics[width=0.92\textwidth]{ch07/fig13_weights}
  \caption{Weights of the ten industries in the frontier portfolio with excess mean $m$, from the
  full-sample moments, \fqFirst{} to \fqLast. The dotted lines mark the means of the global
  minimum-variance and tangency portfolios, and each line is named at its right end. Away from the
  tangency portfolio the short positions and the gross leverage grow quickly: at \fwDoubleMu\% a year
  the portfolio is short \fwDoubleShort\% in \fwDoubleShortName.}
  \label{fig7:weights}
\end{figure}


\paragraph{Adding a riskless asset.}
With a riskless asset paying $r_f$, an investor can combine any risky portfolio with lending or
borrowing, and the mean and volatility of the combination lie on a straight line from $(0,r_f)$
\parencite{tobin1958}. The best line is the steepest, so the question becomes which portfolio has
the highest Sharpe ratio $\theta(w)=w^\top\mu_x/\sqrt{w^\top\Sigma w}$.

\begin{proposition}[The tangency portfolio]\label{prop7:tangency}
Among portfolios with $w^\top\mu_x>0$ the Sharpe ratio is largest for $w\propto\Sigma^{-1}\mu_x$, and
the largest squared Sharpe ratio is
\begin{equation}\label{eq7:theta}
  \theta^2=\mu_x^\top\Sigma^{-1}\mu_x .
\end{equation}
If $\mathbf 1^\top\Sigma^{-1}\mu_x>0$, which holds exactly when $r_f$ is below the mean of the global
minimum-variance portfolio, the fully invested version $w_T=\Sigma^{-1}\mu_x/\mathbf 1^\top\Sigma^{-1}\mu_x$
lies on the efficient branch of the frontier.
\end{proposition}

\begin{proof}
By Cauchy-Schwarz, $w^\top\mu_x=(\Sigma^{1/2}w)^\top(\Sigma^{-1/2}\mu_x)\le\sqrt{w^\top\Sigma w}\sqrt{\mu_x^\top\Sigma^{-1}\mu_x}$,
with equality exactly when $\Sigma^{1/2}w$ is a positive multiple of $\Sigma^{-1/2}\mu_x$. Since
$\mu_x=\mu-r_f\mathbf 1$, $\mathbf 1^\top\Sigma^{-1}\mu_x=B-r_fA$, positive exactly when $r_f<B/A$.
Then $w_T$ is of the form \cref{eq7:frontier} with $\lambda=-r_f/(B-r_fA)$ and $\gamma=1/(B-r_fA)>0$,
and $\gamma>0$ is the condition for the efficient branch, since $\gamma=(Am-B)/D$.
\end{proof}

\Cref{fig7:frontier} shows the frontier of the ten industries, with means and covariances estimated
over the whole sample, in excess returns so that the riskless rate is zero. The global
minimum-variance portfolio has an excess mean of \frGmvMu\% and a volatility of \frGmvVol\%. The
tangency portfolio has a mean of \frTanMu\%, a volatility of \frTanVol\% and a Sharpe ratio of
\frTanSR, against \frMktSR{} for the market factor. To get there it sells short \frShortMax\% of its
capital in \frShortName{} and holds gross positions of \frGross{} times its capital
(\cref{tab7:weights}). These are in-sample numbers: the same data chose the weights and judged them,
a point taken up in \cref{sec7:estimation}.
The listing below computes $A$, $B$ and $C$ from the data files and the two portfolios of
\cref{tab7:weights} from their closed forms; its last line checks that the minimum-variance mean is
$B/A$ and that the tangency Sharpe ratio is $\sqrt{12\,C}$ a year when the riskless rate is zero.
\Needspace*{21\baselineskip}
\lstinputlisting[language=Python]{code/ch07/snippets/closed_form.py}

\begin{figure}[!tb]
  \centering
  \includegraphics[width=0.86\textwidth]{ch07/fig02_frontier}
  \caption{Minimum-variance frontier of the ten industry portfolios in excess returns, \fqFirst{}
  to \fqLast, with short sales allowed and with weights restricted to be nonnegative, the capital
  market line through the tangency portfolio, and the market factor. Only the efficient part of the
long-only frontier is drawn. It ends at
  the industry with the highest mean, because a long-only portfolio cannot have a higher mean than its
  best asset.}
  \label{fig7:frontier}
\end{figure}


\begin{table}[t]
\centering
\caption{Weights in percent of five portfolios of the ten industries, with means and covariances
estimated from \fqFirst{} to \fqLast. Market weights are those of \fqLast. The last three rows use
the full-sample moments for every column.}
\label{tab7:weights}
\small
\begin{tabular}{lrrrrr}
\toprule
 & minimum & tangency & minimum & tangency & market\\
 & variance & & variance, & long only & weights\\
 & & & long only & & \\
\midrule
\frTableRows
\bottomrule
\end{tabular}
\end{table}

\paragraph{Long-only portfolios.}
Most allocators cannot sell short, and many cannot borrow. Adding $w\ge0$ removes the closed form but
keeps the problem convex.

\begin{proposition}[Long-only portfolios]\label{prop7:longonly}
Let $M$ be a matrix with full row rank and $b$ a vector such that some $w\ge0$ satisfies $Mw=b$. The
problem of minimising $\frac12w^\top\Sigma w$ subject to $Mw=b$ and $w\ge0$ has a unique solution,
and $w$ is that solution exactly when there are multipliers $\lambda$ and $\nu$ with
\begin{equation}\label{eq7:kkt}
  \Sigma w=M^\top\lambda+\nu,\qquad \nu\ge0,\qquad \nu_iw_i=0\ \text{for every } i .
\end{equation}
On the set of assets it holds, the solution solves the equality-constrained problem for those assets
alone, and when $M$ stacks $\mathbf 1^\top$ and $\mu^\top$ it is a frontier portfolio of those assets. If some
asset has a positive excess mean, the long-only portfolio with the largest Sharpe ratio is $y/\mathbf 1^\top y$,
where $y$ solves the problem with $M=\mu_x^\top$ and $b=1$.
\end{proposition}

\begin{proof}
The feasible set is a nonempty closed convex polyhedron and the objective is strictly convex and grows
without bound, so a unique minimum exists, and for a convex problem with linear constraints the
Karush-Kuhn-Tucker conditions \cref{eq7:kkt} are necessary and sufficient. Deleting the assets with
$w_i=0$ leaves $\Sigma_{ss}w_s=M_s^\top\lambda$ on the support $s$, the stationarity condition of the
equality-constrained problem on $s$. For the last statement, any $w\ge0$ with $w^\top\mu_x>0$ can be
rescaled to $y=w/w^\top\mu_x$, which satisfies $\mu_x^\top y=1$, and $\theta(w)=\theta(y)=1/\sqrt{y^\top\Sigma y}$.
Maximising the Sharpe ratio is therefore the same as minimising $y^\top\Sigma y$, and rescaling the
minimiser to add to one does not change its Sharpe ratio.
\end{proof}

The second statement of \cref{prop7:longonly} gives an exact algorithm for a handful of assets: for
each candidate support solve the linear equations of the equality-constrained problem, and keep the
candidate whose weights are nonnegative and whose multipliers $\nu_i$ on the assets left out are
nonnegative; the multipliers $\lambda$ of the equality constraints have no sign restriction. With ten industries there are 1,023
supports, and \texttt{ch07lib.qp\_enum} tries them all. The R script runs its own search over supports
and Octave's \texttt{qp} solves the same problems by an active-set method; all three agree. The
long-only tangency portfolio holds \frLoNword{} industries and has a Sharpe ratio of \frLoSR, so in
sample the short positions raise the Sharpe ratio from \frLoSR{} to \frTanSR. The long-only frontier in
\cref{fig7:frontier} lies inside the unconstrained one and touches it only where the unconstrained
weights happen to be nonnegative.
The Octave version uses the scaling in the proof of \cref{prop7:longonly}: minimise
$\frac12y^\top\Sigma y$ subject to $\mu_x^\top y=1$ and $y\ge0$, then divide $y$ by its sum. It prints
the Sharpe ratio of \frLoSR{} and the \frLoNword{} industries held.
\Needspace*{12\baselineskip}
\lstinputlisting[language=Octave]{code/ch07/snippets/lo_tangency.m}

\FloatBarrier
% =====================================================================================
\section{What estimation error does to the frontier}
\label{sec7:estimation}

The frontier needs $\mu$ and $\Sigma$, and both are estimated. Covariances are estimated far more
precisely than means, though with many assets and few years of monthly data the sample covariance
matrix is noisy too and can be close to singular. The means are much harder: with a volatility of 15\% a year, the
standard error of a mean estimated from five years is about $15/\sqrt5\approx6.7$ points, as large as
the means themselves. An optimiser treats estimates as truth and leans hardest on the assets whose
means are most overstated, which led \textcite{michaud1989} to call mean-variance optimisers error
maximisers. The bias in the maximum Sharpe ratio can be computed exactly for normal returns.

\begin{lemma}[Moments of the sample mean and covariance]\label{lem7:wishart}
Let $x_1,\dots,x_T$ be independent $\N(\mu_x,\Sigma)$ vectors of dimension $N$, with sample mean
$\hat\mu$ and sample covariance $\hat\Sigma$ computed with divisor $T-1$. Then $\hat\mu$ and $\hat\Sigma$
are independent, $\hat\mu\sim\N(\mu_x,\Sigma/T)$, and for $T>N+2$
\begin{equation}\label{eq7:wishart}
  \E\hat\Sigma^{-1}=\frac{T-1}{T-N-2}\,\Sigma^{-1} .
\end{equation}
\end{lemma}

This is a standard property of the Wishart distribution, since $(T-1)\hat\Sigma$ is Wishart with
$T-1$ degrees of freedom; \textcite{muirhead1982} proves it.

\begin{proposition}[The in-sample Sharpe ratio is biased upwards]\label{prop7:bias}
Under the assumptions of \cref{lem7:wishart}, the largest squared Sharpe ratio computed from the sample
moments, $\hat\theta^2=\hat\mu^\top\hat\Sigma^{-1}\hat\mu$, has mean
\begin{equation}\label{eq7:bias}
  \E\hat\theta^2=\frac{T-1}{T-N-2}\Bigl(\theta^2+\frac NT\Bigr).
\end{equation}
The true Sharpe ratio of the estimated portfolio $\hat\Sigma^{-1}\hat\mu$ is never above $\theta$.
\end{proposition}

\begin{proof}
By independence, $\E\hat\theta^2=\E\operatorname{tr}(\hat\Sigma^{-1}\hat\mu\hat\mu^\top)=\operatorname{tr}\bigl(\E\hat\Sigma^{-1}\,\E\hat\mu\hat\mu^\top\bigr)$.
By \cref{lem7:wishart} this is $\frac{T-1}{T-N-2}\operatorname{tr}\bigl(\Sigma^{-1}(\Sigma/T+\mu_x\mu_x^\top)\bigr)$,
which is \cref{eq7:bias}. The second statement is \cref{prop7:tangency}: no portfolio has a Sharpe ratio
above $\theta$.
\end{proof}

\begin{figure}[!tb]
  \centering
  \includegraphics[width=0.94\textwidth]{ch07/fig03_estimation}
  \caption{Sharpe ratios of the sample tangency portfolio in \esM{} simulated samples of 60 and 240
  months drawn from normal returns with the full-sample moments of the ten industries. The in-sample
  ratio uses the estimated moments, the true ratio the moments that generated the data. The vertical
  line is the largest Sharpe ratio available, \esTheta.}
  \label{fig7:estimation}
\end{figure}

Two effects push the estimate up. The term $N/T$ is the noise in the means, which the optimiser turns
into apparent Sharpe ratio, and the factor $(T-1)/(T-N-2)$ is the noise in the covariances, which makes
the inverse too large on average. \Cref{fig7:estimation} takes the full-sample moments of the ten
industries as the truth, with $\theta=\esTheta$ a year, and draws \esM{} normal samples of each
length. With five years of data the root mean square of the in-sample Sharpe ratio,
$\sqrt{\E\hat\theta^2}$ in annual units, is \esInSixty, against \esInTheorySixty{} from
\cref{eq7:bias}, which gives $\E\hat\theta^2$ exactly; by Jensen's inequality the mean of $\hat\theta$
itself is lower. The average true Sharpe ratio of the estimated portfolio
is only \esOutSixty{}, and in \esNegSixty\% of samples it is negative. With ten years the numbers are
\esInOneTwenty{} and \esOutOneTwenty, and with twenty years \esInTwoForty{} and \esOutTwoForty. Even
twenty years of data recover two thirds of the available Sharpe ratio while promising half again as
much. \textcite{kan2007} derive the expected out-of-sample loss and show how to shrink the weights
towards the minimum-variance portfolio to reduce it.

\Cref{fig7:convergence} follows both quantities from 36 to \cvTmax{} months, with \cvM{} samples at each
length. The simulated in-sample values lie on the curve from \cref{eq7:bias}, and even at \cvTmax{}
months, fifty years of data, the in-sample ratio is \cvInMax{} while the estimated portfolio delivers
\cvOutMax, \cvShareMax\% of the true \esTheta.

\begin{figure}[!tb]
  \centering
  \includegraphics[width=0.86\textwidth]{ch07/fig14_convergence}
  \caption{Root mean square in-sample Sharpe ratio of the sample tangency portfolio, exact from
  \cref{eq7:bias} and simulated, and the average true Sharpe ratio of the estimated portfolio, against
  the length of the sample, \cvM{} normal samples at each length with the full-sample moments of the
  ten industries. The dashed line is the largest Sharpe ratio available, \esTheta.}
  \label{fig7:convergence}
\end{figure}



\FloatBarrier
% =====================================================================================
\section{The CAPM as a statement about one portfolio}
\label{sec7:capm}

If every investor uses the same means and covariances, cares only about the mean and variance of
the portfolio, and can lend and borrow at $r_f$, then by
\cref{prop7:tangency} every investor holds the same risky portfolio $w_T$ in some proportion. Adding
up across investors, with markets clearing and the riskless asset in zero net supply, the market
portfolio of risky assets must be $w_T$. This is the capital asset
pricing model of \textcite{sharpe1964} and \textcite{lintner1965}. Its content is captured by an
identity that holds for any portfolio, whether or not anyone is in equilibrium.

\begin{proposition}[Beta pricing and efficiency]\label{prop7:capm}
Let $p$ be a portfolio with weights $w_p$, excess mean $\mu_{x,p}\ne0$ and variance $\sigma_p^2>0$, and
let $\beta_i=\Cov(R_i,R_p)/\sigma_p^2$. Then
\begin{equation}\label{eq7:capm}
  \mu_{x,i}=\beta_i\,\mu_{x,p}\qquad\text{for every asset } i
\end{equation}
if and only if $w_p$ is a multiple of $\Sigma^{-1}\mu_x$, that is, if and only if $p$ has the largest
Sharpe ratio or the smallest.
\end{proposition}

\begin{proof}
The vector of betas is $\beta=\Sigma w_p/\sigma_p^2$. If $w_p=c\Sigma^{-1}\mu_x$, then $\Sigma w_p=c\mu_x$
and $\sigma_p^2=w_p^\top\Sigma w_p=c\,\mu_{x,p}$, so $\beta=\mu_x/\mu_{x,p}$, which is \cref{eq7:capm}.
Conversely, \cref{eq7:capm} says $\mu_x=\mu_{x,p}\Sigma w_p/\sigma_p^2$, so
$w_p=(\sigma_p^2/\mu_{x,p})\Sigma^{-1}\mu_x$. The sign of the multiple decides whether $p$ maximises or
minimises the Sharpe ratio, by \cref{prop7:tangency}.
\end{proof}

The line $\mu_{x,i}=\beta_i\mu_{x,M}$ in the plane of beta and mean is the security market line. Read
one way, \cref{prop7:capm} is a model of expected returns: assets are rewarded only for their
covariance with the market, and an asset with a beta of zero earns the riskless rate however volatile
it is. Read the other way, it is a warning that \textcite{roll1977} made precise: testing the CAPM is
the same as testing whether a market proxy is mean-variance efficient, and a rejection could mean the
proxy is wrong. In the industry data the identity can be checked against the in-sample tangency
portfolio, where it must hold, and it does to rounding error: the right panel of \cref{fig7:sml}
puts all ten industries on the line, with tangency betas from \cpTanBetaMin{} to \cpTanBetaMax. Against
the market factor it need not hold. \textcite{black1972} showed that without riskless borrowing the line still holds with
the mean of a zero-beta portfolio in place of $r_f$, so its intercept can be above the bill rate.

\begin{figure}[!tb]
  \centering
  \includegraphics[width=0.76\textwidth]{ch07/fig04_sml}
  \caption{Mean excess return of each industry against its beta on the market factor (left) and on
  the in-sample tangency portfolio (right), \fqFirst{} to \fqLast. On the left the dashed line is the
  security market line of the CAPM, through the origin with slope equal to the market premium, and the
  solid line is the least-squares fit, flatter and with a positive intercept. On the right the dashed
  line has slope equal to the mean of the tangency portfolio, and the points lie on it exactly.}
  \label{fig7:sml}
\end{figure}

\Cref{fig7:sml} plots the mean excess returns of the industries from \fqFirst{} to \fqLast{} against
their betas on the market. Betas run from \smBetaMin{} for \smBetaMinName{} to \smBetaMax{} for
\smBetaMaxName. The least-squares line has a slope of \smSlope\% a year per unit of beta, with a t
statistic of \smSlopeT, and an intercept of \smIcpt\%, while the CAPM predicts a slope equal to the
market premium of \frMktMu\% and an intercept of zero. The line is too flat: low-beta industries
earned more than the model predicts and high-beta industries less, which is the pattern the zero-beta
version of the model can accommodate. With ten points the evidence is thin, and
\cref{sec7:alpha} turns it into a formal test.


\FloatBarrier
% =====================================================================================
\section{Factor models}
\label{sec7:factors}

A covariance matrix of $N$ assets has $N(N+1)/2$ free entries, 125,250 for five hundred stocks, and
with fewer months than assets the sample matrix is singular. Factor models impose structure.

\begin{definition}[Linear factor model]\label{def7:factor}
Excess returns follow a linear factor model with $K$ factors if
\begin{equation}\label{eq7:factor}
  x_t=\alpha+Bf_t+\varepsilon_t ,
\end{equation}
where $f_t$ has mean $\bar f$ and covariance $\Omega$, $B$ is $N\times K$, the residuals $\varepsilon_t$
have mean zero, are uncorrelated with $f_t$ and have a diagonal covariance matrix $D$. Then
$\Sigma=B\Omega B^\top+D$.
\end{definition}

\begin{proposition}[Inverting a factor covariance]\label{prop7:woodbury}
Under \cref{def7:factor}, with $\Omega$ nonsingular and every residual variance positive,
\begin{equation}\label{eq7:woodbury}
  \Sigma^{-1}=D^{-1}-D^{-1}B\bigl(\Omega^{-1}+B^\top D^{-1}B\bigr)^{-1}B^\top D^{-1},
\end{equation}
so inverting $\Sigma$ needs only the inverse of a $K\times K$ matrix, and the model has
$NK+K(K+1)/2+N$ parameters.
\end{proposition}

\begin{proof}
Write $G=\Omega^{-1}+B^\top D^{-1}B$. Multiplying the right side of \cref{eq7:woodbury} by
$B\Omega B^\top+D$ gives
$I+D^{-1}B\Omega B^\top-D^{-1}BG^{-1}\bigl(B^\top D^{-1}B\Omega+I\bigr)B^\top$. Since
$B^\top D^{-1}B\Omega+I=G\Omega$, the last term is $D^{-1}B\Omega B^\top$, which cancels the second.
The count adds the loadings, the free entries of $\Omega$ and the diagonal of $D$.
\end{proof}

With five hundred stocks and six factors the count is 3,521 instead of 125,250. With ten industries a
six-factor model has more parameters than the sample covariance, so its value here is as a
description. Fitted with the factors of \textcite{fama2015} and the momentum factor of
\textcite{carhart1997}, the model reproduces the industry correlations to within \fmCorrErr{}, and its
minimum-variance portfolio has a predicted volatility of \fmGmvFactor\% against \fmGmvFactorTrue\%
when evaluated with the sample covariance; the diagonal $D$ ignores the residual correlation between
industries, so the model is too optimistic about how much can be hedged. \Cref{fig7:corrgap} shows
where the fit fails. The gaps have both signs, positive for \fmGapPos{} of the \fmGapPairs{} pairs, and
the largest, \fmGapMax, is for \fmGapPair, whose residuals move against each other once the factors
are removed.

\begin{figure}[!tb]
  \centering
  \includegraphics[width=0.74\textwidth]{ch07/fig15_corrgap}
  \caption{Sample correlation of the ten industries minus the correlation implied by the six-factor
  model $B\Omega B^\top+D$, \fqFirst{} to \fqLast. Each pair appears once; every cell carries its
  sign, with a third decimal where the gap rounds to zero at two.}
  \label{fig7:corrgap}
\end{figure}


\paragraph{Arbitrage pricing.}
A factor model is a statement about covariances. \textcite{ross1976} showed that it becomes a
statement about means once arbitrage is ruled out. The exact version has a short proof.

\begin{proposition}[Exact arbitrage pricing]\label{prop7:apt}
Suppose $x_t=\mu_x+B(f_t-\bar f)$ with no residual, and that every portfolio $w$ with $B^\top w=0$ has
$w^\top\mu_x=0$. Then $\mu_x=B\lambda$ for some $\lambda\in\mathbb R^K$. If the factors are themselves
excess returns of traded portfolios and these are included among the assets, then $\lambda=\bar f$ and $\alpha=0$ in \cref{eq7:factor}.
\end{proposition}

\begin{proof}
A portfolio with $B^\top w=0$ has excess return $w^\top\mu_x$ with no risk, which must be zero, so
$\mu_x$ is orthogonal to the null space of $B^\top$. That space is the orthogonal complement of the
column space of $B$, so $\mu_x$ lies in the column space. If factor $k$ is the excess return of a
traded portfolio, include it among the assets: its loadings are the $k$th unit vector and its mean is
$\bar f_k$, so $\lambda_k=\bar f_k$. Then $\alpha=\mu_x-B\bar f=0$.
\end{proof}

With residual risk the conclusion holds only approximately. \textcite{huberman1982} shows that if
the returns have a factor structure, the residual variances are bounded by a common constant and no
sequence of portfolios gives an arbitrage in the limit, the sum of squared pricing errors across assets
stays bounded as $N$ grows, so most assets are priced nearly right. The bound is on the errors taken
together and is a statement about large $N$; it does not say that any one asset has a small alpha. This is why the
intercept of \cref{eq7:factor} with traded factors is the natural definition of alpha: it is what a
manager earns beyond a portfolio of the factors with the same loadings.

\begin{figure}[!tb]
  \centering
  \includegraphics[width=0.94\textwidth]{ch07/fig05_pca}
  \caption{Principal components of the correlation matrix of the ten industries, \fqFirst{} to
  \fqLast. Left, the share of total variance explained by each component. Right, the loadings of the
  first two; the second is signed so that technology has a positive loading.}
  \label{fig7:pca}
\end{figure}

\paragraph{Factors from the data.}
The factors need not be chosen in advance. The eigenvectors of the correlation matrix of the
industries are weights on the standardised returns; the combinations they form are uncorrelated, and
their eigenvalues are the variances they explain. Combinations of the raw returns with the same
weights are not uncorrelated, and the weights need not add to one, so a component is not by itself a
fully invested portfolio. \Cref{fig7:pca} shows that the first component explains \pcOne\% of the variance of the
standardised returns and the first three \pcThree\%. The first has positive loadings on every industry
and a correlation of \pcMktCorr{} with the market factor. The second, with \pcTwo\%, sets technology,
durable goods and retail against utilities and energy; its correlation with the value factor is
\pcTwoHmlSigned, so it leans towards growth against value without being a value portfolio. Statistical factors explain the
covariance well, but they come with no theory of their means, which is where the characteristic
factors of \textcite{fama1993} come in.


\FloatBarrier
% =====================================================================================
\section{Estimating alpha and beta}
\label{sec7:alpha}

A manager's alpha and beta are estimated by the time-series regression
\begin{equation}\label{eq7:reg}
  x_{p,t}=\alpha+\beta^\top f_t+\varepsilon_t,\qquad t=1,\dots,T,
\end{equation}
of the fund's excess return on the excess returns of $K$ traded factors, the single market factor in
the CAPM. When the traded factors price the assets, \cref{prop7:apt} gives $\alpha=0$;
in general $\alpha$ is the mean return left after paying for the factor exposures at the factors' own
mean returns, and reading it as skill assumes the factors are the right benchmark. Least squares gives the estimates, and its classical standard
errors assume residuals that are uncorrelated over time with constant variance. Monthly fund returns
often fail both assumptions, and the errors of \textcite{newey1987}, used throughout this section with
\alNW{} lags, allow for both. \Cref{tab7:alphas} reports the regressions for the ten industries.

\begin{table}[t]
\centering
\caption{Alphas, in percent a year, of the ten industries, \fqFirst{} to \fqLast, against the market
factor alone (CAPM), with the beta and $R^2$, and against the market, size, value, profitability,
investment and momentum factors (six factors), with the $R^2$. The t statistics use Newey-West standard errors with \alNW{} lags.}
\label{tab7:alphas}
\small
\begin{tabular}{lrrrrrrr}
\toprule
 & \multicolumn{4}{c}{CAPM} & \multicolumn{3}{c}{six factors}\\
\cmidrule(lr){2-5}\cmidrule(lr){6-8}
 & $\alpha$ & $t$ & $\beta$ & $R^2$ & $\alpha$ & $t$ & $R^2$\\
\midrule
\alTableRows
\bottomrule
\end{tabular}
\end{table}

Against the market alone none of the ten alphas is significant at the 5\% level, and the market
explains between \alRsqMinCapm{} and \alRsqMaxCapm{} of the variance. Adding five factors raises the
explained variance to between \alRsqMinSix{} and \alRsqMaxSix, and yet \alNsigSixWord{} industries now
have significant alphas, among them technology at \alHiTecA\% a year with a t statistic of
\alHiTecT{} and the residual industry at \alOtherA\%. More factors explain more of the variance and
can make the pricing errors worse: technology loads negatively on the value, profitability,
investment and momentum factors, whose premia were all positive, so the six-factor model predicts a lower mean for it
than the CAPM does. \Cref{sec7:attribution} breaks this down.

The listing below runs the two regressions for health care with statsmodels and prints both kinds of
t statistic. The Newey-West option needs \texttt{use\_correction=True} to apply the small-sample
factor $T/(T-k)$ used in \cref{tab7:alphas}; without it the t statistics come out slightly larger.
Classical and Newey-West t statistics differ little for an industry portfolio, but for a fund
with smoothed or option-like returns the gap can be large.
\Needspace*{18\baselineskip}
\lstinputlisting[language=Python]{code/ch07/snippets/hac_alpha.py}

\paragraph{How precise is an alpha?}
The standard error of the intercept has a simple form that explains why alphas are so hard to pin
down.

\begin{proposition}[Standard error of alpha]\label{prop7:alphase}
In the regression $y_t=\alpha+\beta x_t+\varepsilon_t$ with residuals independent of $x$ and of each
other, with mean zero and variance $\sigma_\varepsilon^2$, the least-squares intercept has
\begin{equation}\label{eq7:alphase}
  \Var(\hat\alpha\mid x)=\frac{\sigma_\varepsilon^2}{T}\Bigl(1+\frac{\bar x^2}{s_x^2}\Bigr),
  \qquad s_x^2=\frac1T\sum_{t=1}^T(x_t-\bar x)^2 .
\end{equation}
\end{proposition}

\begin{proof}
The estimates are $\hat\beta=\sum_t(x_t-\bar x)y_t/(Ts_x^2)$ and $\hat\alpha=\bar y-\hat\beta\bar x$.
Given $x$, $\Var(\bar y)=\sigma_\varepsilon^2/T$, $\Var(\hat\beta)=\sigma_\varepsilon^2/(Ts_x^2)$ and
$\Cov(\bar y,\hat\beta)=\sigma_\varepsilon^2\sum_t(x_t-\bar x)/(T^2s_x^2)=0$. Hence
$\Var(\hat\alpha)=\sigma_\varepsilon^2/T+\bar x^2\sigma_\varepsilon^2/(Ts_x^2)$.
\end{proof}

The ratio $\bar x/s_x$ is the sample Sharpe ratio of the factor, which is small for monthly data, so
the standard error of alpha is essentially that of a mean of the residuals. For health care the
residual volatility is \alHlthResid\% a year, the CAPM alpha \alHlthA\% and its classical standard
error \alHlthSe\%, which \cref{eq7:alphase} reproduces exactly. The ratio of alpha to residual
volatility is the information ratio, here \alHlthIR, and the ratio decides how long a track record must
be.

\begin{corollary}[Years needed to detect alpha]\label{cor7:detect}
With information ratio $\mathrm{IR}=\alpha/\sigma_\varepsilon$ and factor Sharpe ratio $\theta_f$, both per
period, stationary factor moments and residuals that are independent with constant variance, the
t statistic of $\hat\alpha$ is centred close to $\mathrm{IR}\sqrt T/\sqrt{1+\theta_f^2}$, and its centre
reaches $t^*$ after
\begin{equation}\label{eq7:detect}
  T^*=\frac{t^{*2}\,(1+\theta_f^2)}{\mathrm{IR}^2}\quad\text{periods.}
\end{equation}
\end{corollary}

With monthly data, an annual information ratio of $\mathrm{IR}_a$ is $\mathrm{IR}_a/\sqrt{12}$ a month,
so $T^*$ in years is close to $4/\mathrm{IR}_a^2$ for $t^*=2$. With the market's Sharpe ratio of
\dtSRm{} a year, a manager with an annual information ratio of 0.5, which is good, needs \dtYrsHalf{}
years for the expected t statistic to reach two; one with 0.25 needs \dtYrsQuarter{} years, and only an
exceptional ratio of one brings it down to \dtYrsOne. The left panel of \cref{fig7:detect} shows the
probability that the t statistic exceeds 1.96 in absolute value. These are the horizons at which the centre of the
statistic reaches two. They guarantee nothing: at that point a test at 5\% rejects only about half the time.

\begin{figure}[!tb]
  \centering
  \includegraphics[width=0.96\textwidth]{ch07/fig06_detect}
  \caption{Left, the probability that the t statistic of alpha exceeds 1.96 in absolute value, against
  the length of the monthly record, for three annual information ratios, from \cref{cor7:detect}.
  Right, t statistics of \mgN{} simulated managers over \mgT{} months, \mgSkill{} of them with an
  information ratio of 0.5 and the rest with none; the dashed line is the 5\% critical value.}
  \label{fig7:detect}
\end{figure}


\paragraph{Testing many alphas at once.}
Ten t statistics are not ten independent tests: the industries' residuals are correlated. The test of
\textcite{gibbons1989} asks whether all $N$ intercepts are zero together. With residual covariance
$\hat\Sigma_\varepsilon$ and factor covariance $\hat\Omega$ estimated with divisor $T$, and factor means $\bar f$,
\begin{equation}\label{eq7:grs}
  J=\frac{T-N-K}{N}\,\frac{\hat\alpha^\top\hat\Sigma_\varepsilon^{-1}\hat\alpha}{1+\bar f^\top\hat\Omega^{-1}\bar f}
\end{equation}
has an $F$ distribution with $N$ and $T-N-K$ degrees of freedom when the intercepts are zero and the
residuals are normal and independent over time, independent of the factors, with a nonsingular
covariance matrix and $T>N+K$. The quadratic form in the numerator has a meaning of
its own.

\begin{proposition}[Alphas as a gain in Sharpe ratio]\label{prop7:grs}
Let $z=(f,x)$ stack the factor and asset excess returns, with $x=\alpha+Bf+\varepsilon$, $\Cov(f)=\Omega$
and $\Cov(\varepsilon)=\Sigma_\varepsilon$. Then
\begin{equation}\label{eq7:grsid}
  \alpha^\top\Sigma_\varepsilon^{-1}\alpha=\theta^2(f,x)-\theta^2(f),
\end{equation}
where $\theta^2(\cdot)$ is the largest squared Sharpe ratio of \cref{eq7:theta} for the assets listed.
The same identity holds for the sample moments with divisor $T$.
\end{proposition}

\begin{proof}
The mean of $z$ is $m=(\bar f,\alpha+B\bar f)$ and its covariance is $V=\begin{psmallmatrix}\Omega&\Omega B^\top\\ B\Omega&B\Omega B^\top+\Sigma_\varepsilon\end{psmallmatrix}$,
whose inverse is $\begin{psmallmatrix}\Omega^{-1}+B^\top\Sigma_\varepsilon^{-1}B&-B^\top\Sigma_\varepsilon^{-1}\\-\Sigma_\varepsilon^{-1}B&\Sigma_\varepsilon^{-1}\end{psmallmatrix}$,
as multiplying out confirms. Expanding $m^\top V^{-1}m$, every term in $B\bar f$ cancels and what is
left is $\bar f^\top\Omega^{-1}\bar f+\alpha^\top\Sigma_\varepsilon^{-1}\alpha$. With sample moments, least squares
makes the residuals orthogonal to the factors and the same algebra applies.
\end{proof}

A joint test of alphas is a test of whether the assets improve the best Sharpe ratio available from
the factors. For the CAPM, $J=\grsCapm$ with a p-value of \grsPCapm: the ten industries raise the
Sharpe ratio from \grsSRcapm{} for the market alone to \grsSRcapmAll{} in sample, and that gain is
within what chance would give. For six factors, $J=\grsSix$ with a p-value of \grsPSix: the factors
alone reach \grsSRfac, and adding the industries reaches \grsSRall. With five factors and no momentum
the p-value is \grsPFive. For these portfolios and dates the joint test rejects zero alphas for the
multifactor models and not for the CAPM. That does not show the CAPM to be the better model in general,
and industries were not the portfolios the factors were built to price.
In R the test takes a few lines: one call to \texttt{lm} with a matrix response, the residual
covariance with divisor $T$, and the $F$ distribution that follows \cref{eq7:grs}.
\Needspace*{19\baselineskip}
\lstinputlisting[language=R]{code/ch07/snippets/grs.R}

\paragraph{Betas move.}
\Cref{fig7:rolling} estimates betas over rolling windows of \rbW{} months. The beta of technology ran
from \rbMinHiTec{} to \rbMaxHiTec, the peak in the window ending in \rbMaxWhenHiTec{} that contains the
collapse of technology stocks after 2000. The beta of utilities fell to \rbMinUtils{} in the window
ending in \rbMinWhenUtils, and that of energy reached \rbMaxEnrgy{} in the window ending in
\rbMaxWhenEnrgy{} before falling to \rbLastEnrgy. The median standard error of a rolling technology
beta is \rbSeHiTec, yet \rbOutside\% of the rolling estimates lie more than two standard errors from
the full-sample beta. Constant beta is a convenience: a fund's beta measured over one window is a
weak guide to its beta in the next.

\begin{figure}[!tb]
  \centering
  \includegraphics[width=0.72\textwidth]{ch07/fig07_rolling}
  \caption{Market betas of three industries estimated over rolling windows of \rbW{} months, plotted at
  the end of each window, with bands of two classical standard errors.}
  \label{fig7:rolling}
\end{figure}


\paragraph{Screening many managers.}
An allocator who screens a database of managers by their t statistics faces the problem in reverse:
with many managers and no skill, some will look skilled. The right panel of \cref{fig7:detect}
simulates \mgN{} managers with ten years of monthly returns, betas between 0.6 and 1.4 and a tracking
error of 8\% a year. One in ten has an annual information ratio of 0.5 and the rest have none. A
two-sided test at 5\% flags \mgNaive{} managers, and \mgNaiveFalse{} of them have no skill; of the
\mgNaivePos{} flagged with positive alphas, \mgNaivePosFalse{} are lucky. By \cref{cor7:detect} a
skilled manager is detected with probability of only \mgPowerTheory\% in ten years. The rule of
\textcite{benjamini1995} controls the expected share of false discoveries among those flagged: sort
the $m$ p-values, find the largest $k$ with $p_{(k)}\le qk/m$, and flag the $k$ smallest. With
$q=10\%$ it flags none of the \mgN{} managers, because the skilled ones are not separated from the
lucky ones. With twenty-year records it flags \mgLongBh, of whom \mgLongBhFalse{} are false, while the
naive test flags \mgLongNaive{} with \mgLongNaiveFalse{} false. Twice the share of p-values above one half
estimates the share of managers without skill \parencite{storey2002}; with ten-year records it gives \mgPiZero, against
a true value of 0.90. \textcite{barras2010} apply these ideas to mutual funds and
\textcite{harvey2016} to the hundreds of factors proposed in the literature, where the same
arithmetic applies to researchers instead of managers.
The listing applies the rule with R's \texttt{p.adjust} to the p-values of the simulated managers,
which the Python script writes out, and prints the counts above with Storey's estimate.
\Needspace*{12\baselineskip}
\lstinputlisting[language=R]{code/ch07/snippets/bh.R}

\FloatBarrier
% =====================================================================================
\section{Stale prices: beta hidden, alpha manufactured}
\label{sec7:smoothing}

Private equity, real estate, infrastructure and many credit and hedge funds report returns based on
appraisals or on prices of assets that rarely trade. Chapter~1 showed that such returns are smoothed:
the reported return mixes this period's true return with earlier ones. A regression of a smoothed
series on the contemporaneous market return then measures the wrong thing.

\begin{proposition}[Regressions on smoothed returns]\label{prop7:stale}
Let the true excess return be $r_t=a+\beta m_t+e_t$, where the market excess returns $m_t$ are
independent and identically distributed with mean $\bar m$ and variance $\sigma_m^2$, and $e_t$ is
white noise independent of the market with variance $\sigma_e^2$. Let the reported return be
$r^o_t=\sum_{k\ge0}\vartheta_kr_{t-k}$, with nonnegative weights $\vartheta_k$ that add to one. Then
\begin{enumerate}[label=(\roman*)]
  \item $\E r^o_t=\E r_t$;
  \item the population slope of $r^o_t$ on $m_t$ alone is $\vartheta_0\beta$, and its intercept is
  $a+(1-\vartheta_0)\beta\bar m$;
  \item in the regression of $r^o_t$ on $m_t,m_{t-1},\dots,m_{t-q}$ the slopes add to
  $\beta\sum_{j=0}^q\vartheta_j$, which tends to $\beta$ as $q$ grows;
  \item $\Var r^o_t=(\beta^2\sigma_m^2+\sigma_e^2)\sum_k\vartheta_k^2\le\Var r_t$.
\end{enumerate}
\end{proposition}

\begin{proof}
(i) holds because the weights add to one. Since the market returns are independent,
$\Cov(r^o_t,m_{t-j})=\vartheta_j\beta\sigma_m^2$, and the lagged market returns are uncorrelated with
each other. The population regression on $m_t,\dots,m_{t-q}$ therefore has the separate slopes
$\vartheta_j\beta$, which gives (ii) and (iii); the intercept in (ii) is the mean of $r^o_t$ less
$\vartheta_0\beta\bar m$. For (iv), $r^o_t$ is a weighted sum of independent terms $\beta m_{t-k}+e_{t-k}$,
and $\sum_k\vartheta_k^2\le\bigl(\sum_k\vartheta_k\bigr)^2=1$.
\end{proof}

The common model of appraisal smoothing is $r^o_t=(1-\phi)r_t+\phi r^o_{t-1}$, for which
$\vartheta_k=(1-\phi)\phi^k$, the slope on the current market is $(1-\phi)\beta$, the sum of $q+1$ slopes
is $(1-\phi^{q+1})\beta$, and $\sum_k\vartheta_k^2=(1-\phi)/(1+\phi)$. The volatility falls by the
square root of that factor while the mean is unchanged, so the Sharpe ratio rises by the factor
$\sqrt{(1+\phi)/(1-\phi)}$, which is $\sqrt3$ for $\phi=0.5$, and
the missing beta times the market premium reappears as alpha. The regression on lagged market
returns is the estimator of \textcite{dimson1979}, proposed for thinly traded shares, and
\textcite{getmansky2004} give a model of the smoothing for hedge funds.

To see the sizes, take the true returns to be the residual industry, whose excess returns have a beta
of \smoBeta{}, a CAPM alpha of \smoAlpha\% a year, a volatility of \smoVol\% and a Sharpe ratio of
\smoSR. The market's own first-order autocorrelation is \smoMktAcf, close enough to zero for
\cref{prop7:stale} to apply. Smoothed with $\phi=0.5$, the series has a first-order autocorrelation of
\smoAcfHalf, a beta of \smoBetaHalf, a volatility of \smoVolHalf\% and a Sharpe ratio of \smoSRHalf.
For uncorrelated true returns the factor $\sqrt{(1-\phi)/(1+\phi)}$ would give a volatility of
\smoVolFormula\%, a Sharpe ratio of \smoSRFormula{} and an autocorrelation of exactly $\phi$. The
industry's own returns have a first-order autocorrelation of \smoOwnAcf, which raises both the
volatility and the autocorrelation of the smoothed series, and the smoothed mean, started at the first
return, is \smoMeanKept\% of the true mean; together they give \smoSRHalf. Its CAPM alpha is \smoAlphaHalf\% with a t statistic of \smoAlphaTHalf. With $\phi=0.8$ the beta
drops to \smoBetaEight{} and the alpha rises to \smoAlphaEight\% with a t statistic of
\smoAlphaTEight. Nothing has changed except the way the returns were reported. Summing slopes over
\smoQ{} lags gives betas of \smoDimsonHalf{} and \smoDimsonEight{} and alphas of
\smoDimsonAlphaHalf\% and \smoDimsonAlphaEight\%, much closer to the truth of \smoBeta{} and \smoAlpha\%,
though with $\phi=0.8$ twelve lags still leave \smoSpurEight{} points of spurious alpha. \Cref{fig7:smoothing} shows how
many lags are needed: with $\phi=0.8$ six lags give a summed beta of only \smoDimsonSixEight, about four fifths of the true
\smoBeta.

\begin{figure}[!tb]
  \centering
  \includegraphics[width=0.84\textwidth]{ch07/fig08_smoothing}
  \caption{Sum of the slopes of a regression of smoothed returns on the current and $q$ lagged market
  returns, for the residual industry smoothed with two values of $\phi$, \fqFirst{} to \fqLast. Dashed
  lines are $(1-\phi^{q+1})\beta$ from \cref{prop7:stale}; the grey line is the beta of the unsmoothed
  returns.}
  \label{fig7:smoothing}
\end{figure}

Unsmoothing with a known $\phi$ recovers the true series exactly, but $\phi$ has to be estimated, and an
allocator should treat a high Sharpe ratio with high autocorrelation as a reporting effect until shown
otherwise. \Cref{fig7:appraisal} shows what the investor would have seen from \smoDdFirst{} to
\smoDdLast. The true series lost \smoDdTrue\% from peak to trough; smoothed with $\phi=0.5$ the
reported drawdown is \smoDdHalf\%, and with $\phi=0.8$ it is \smoDdEight\%, reached months later.

\begin{figure}[!tb]
  \centering
  \includegraphics[width=0.86\textwidth]{ch07/fig16_appraisal}
  \caption{Wealth from the excess returns of the residual industry, \smoDdFirst{} to \smoDdLast,
  true and smoothed with two values of $\phi$, with the largest peak-to-trough loss of each. The
  smoothing runs over the whole sample.}
  \label{fig7:appraisal}
\end{figure}





\FloatBarrier
% =====================================================================================
\section{Option-like payoffs}
\label{sec7:options}

Many hedge fund strategies have payoffs that are not linear in the market. Merger arbitrage loses
when deals break in a falling market, and a fund that sells index puts earns a steady premium and
loses heavily in crashes; \textcite{agarwal2004} find such exposures across hedge fund styles. A
linear regression assigns these funds a beta and an alpha, and the alpha can have either sign
depending on what the options were worth.

Take a fund that at the start of each month sells one month puts on the market struck at \poMoney\%
of its level, holds the strike in bills as collateral, and prices the puts with the Black-Scholes
formula of Chapter~5 at the trailing twelve month realised volatility. The market path is the actual
one from \poFirst{} to \fqLast; the first year of the data is used for the first volatility estimate.
The pricing rule is hypothetical and is not a record of premiums that could have been traded: it
ignores bid-ask spreads, margin and the gap between implied and realised volatility. Priced at realised volatility, the premium is \poZeroPrem\% a year, the
mean excess return \poZeroMean\% and the CAPM alpha \poZeroA\% with a t statistic of \poZeroT. Priced at a
volatility 20\% higher, an assumption made only to show the effect, the premium is \poTwentyPrem\%, the mean \poTwentyMean\% and the alpha \poTwentyA\% with a t
statistic of \poTwentyT. In both cases the beta is \poTwentyB{} and the worst month is
\poTwentyWorstWhen, with losses of \poZeroWorst\% and \poTwentyWorst\%. The fund's alpha is a statement about the price of
puts, which the CAPM does not see.

\Cref{fig7:puts} shows the shape. \textcite{treynor1966} add the squared market return to the
regression, and \textcite{henriksson1981} add $\max(-m_t,0)$, which with returns measured in excess of bills
is the payoff, per unit of the market, of a one-period put struck at the bill-grown level; the regressor is built from the market return and is
not itself an observed option return. In the second regression, with the 20\% markup, the coefficient on the put
payoff is \poTwentyHmG{} and the $R^2$ rises from \poTwentyCapmR{} to \poTwentyHmR. The intercept,
\poTwentyHmA\% a year, is the fitted return in a flat month, and it sits well above both the \poTwentyPrem\%
average premium the fund collects and the \poTwentyFlat\% it actually earned in the \poTwentyFlatN{}
months when the market moved by less than half a percent: the regressor kinks at zero while the fund's payoff kinks at the
strike, \poMoney\% of the level, so the fit at zero lies above the cloud of flat months in
\cref{fig7:puts}. It is not skill either. What is missing is the value of the put exposure.

\begin{figure}[!tb]
  \centering
  \includegraphics[width=0.84\textwidth]{ch07/fig09_puts}
  \caption{Monthly excess returns of a fund that sells one month puts struck at \poMoney\% of the
  market, priced at 1.2 times trailing realised volatility, against the market excess return,
  \poFirst{} to \fqLast, with the CAPM line and the Henriksson-Merton fit.}
  \label{fig7:puts}
\end{figure}

\begin{proposition}[Alpha against traded options]\label{prop7:gj}
Suppose $x_{p,t}=a+b\,m_t+c\,h_t+\varepsilon_t$, where $h_t$ is the payoff of an option on the market
that costs $\pi_t$ at the start of the period, both measured as fractions of the market's level at the start and
with the cost carried to the end of the period at the bill rate, and $\varepsilon_t$ has mean zero. A
self-financing portfolio of $b$ in the market and $c$ options, with the balance in bills, has excess return $b\,m_t+c(h_t-\pi_t)$, and the fund's
excess return over that portfolio has mean $a+c\,\E\pi_t$.
\end{proposition}

\begin{proof}
An option bought for $\pi_t$ with borrowed money pays $h_t-\pi_t$ in excess of the bill. Subtracting
the replicating portfolio from the fund leaves $a+c\pi_t+\varepsilon_t$, whose mean is $a+c\,\E\pi_t$.
\end{proof}

This is the idea of \textcite{glosten1994}: measure performance against a benchmark that includes the
options the fund could have traded. With the at-the-money put priced at the same volatility as the
fund's own, the alpha of \cref{prop7:gj} is \poZeroAdj\% a year for the fund that sells at realised
volatility and \poTwentyAdj\% for the one that sells at the markup. The CAPM alphas of \poZeroA\% and
\poTwentyA\% were swings in the price of insurance. Against traded options both funds earn slightly less than
they would by selling the at-the-money put of the benchmark, because a put struck at \poMoney\% is a
different contract. Because the benchmark put is priced at the fund's own volatility, this alpha compares two strikes, not
two prices of insurance, so the markup moves this alpha by only \poMarkupAdj{} points, against
\poMarkupCapm{} for the CAPM alpha. For the fund priced at the
markup, the squared-market regression tells the same story through a coefficient of \poTwentyTmG{}
with a t statistic of \poTwentyTmT.


\FloatBarrier
% =====================================================================================
\section{Return attribution}
\label{sec7:attribution}

Attribution asks where a portfolio's return relative to its benchmark came from. A factor
attribution uses \cref{eq7:reg}: the mean excess return is $\hat\alpha+\sum_k\hat\beta_k\bar f_k$
exactly, because least squares fits the means. For the technology industry over \fqFirst{} to \fqLast,
the mean excess return of \faMean\% a year splits into \faMkt\% from the market, \faSmb\% from size,
\faHml\% from value, \faRmw\% from profitability, \faCma\% from investment, \faMom\% from momentum and
an alpha of \faAlpha\% (the components are rounded: the printed ones add to \faRoundedSum, and the four negative factor terms
to \faFourRoundedNeg, against \faFourNegSigned{} unrounded). Against the CAPM's predicted mean of \faCapmPred\%, the smaller market loading takes off \faMktDrop{}
points and the negative loadings on the last four factors \faFourNeg{} more, offset by \faSmb{} from
size; that gap of \faGap{} points is what turns a CAPM alpha of \alHiTecCapm\% into a six-factor alpha
of \faAlpha\% in \cref{tab7:alphas}. \Cref{fig7:waterfall} draws the same split.

\begin{figure}[!tb]
  \centering
  \includegraphics[width=0.86\textwidth]{ch07/fig17_waterfall}
  \caption{Mean excess return of technology, \fqFirst{} to \fqLast, split into the contributions
  $\hat\beta_k\bar f_k$ of the six factors and the alpha, in percent a year. The bars are computed from
  unrounded estimates and add to the mean exactly. The dashed line is the mean predicted by the CAPM.}
  \label{fig7:waterfall}
\end{figure}


A holdings-based attribution splits the active return by segments, here industries. Let the
portfolio hold weight $w_i$ in segment $i$ and earn $R^p_i$ there, the benchmark hold $b_i$ and earn
$R^b_i$, and $R_b=\sum_ib_iR^b_i$.

\begin{proposition}[Allocation, selection and interaction]\label{prop7:brinson}
The active return $R_p-R_b$ of one period is the sum over segments of
\begin{equation}\label{eq7:brinson}
  \underbrace{(w_i-b_i)(R^b_i-R_b)}_{\text{allocation}}+\underbrace{b_i(R^p_i-R^b_i)}_{\text{selection}}
  +\underbrace{(w_i-b_i)(R^p_i-R^b_i)}_{\text{interaction}} .
\end{equation}
\end{proposition}

\begin{proof}
The three terms add to $w_iR^p_i-b_iR^b_i-(w_i-b_i)R_b$, and the last term sums to zero over segments
because both sets of weights add to one.
\end{proof}

The decomposition is that of \textcite{brinson1986}, with the allocation effect measured against the
benchmark return as in \textcite{brinson1985}, so that overweighting a segment that beat the benchmark
counts as good allocation. Over several periods the effects do not add up: active returns compound,
and the sum of single-period active returns is not the active return of the whole period. Linking
factors fix this.

\begin{proposition}[Linking over periods]\label{prop7:carino}
Let $R_t$ and $B_t$ be portfolio and benchmark returns in periods $t=1,\dots,n$, $R$ and $B$ the
compounded returns, and
\[
  k_t=\frac{\ln(1+R_t)-\ln(1+B_t)}{R_t-B_t},\qquad k=\frac{\ln(1+R)-\ln(1+B)}{R-B},
\]
with $k_t=1/(1+R_t)$ when $R_t=B_t$ and $k=1/(1+R)$ when $R=B$, all gross returns $1+R_t$, $1+B_t$
being positive. Then $\sum_t(k_t/k)(R_t-B_t)=R-B$, so the effects of
\cref{eq7:brinson} scaled by $k_t/k$ add up over periods to the compounded active return.
\end{proposition}

\begin{proof}
$\sum_tk_t(R_t-B_t)=\sum_t\ln(1+R_t)-\sum_t\ln(1+B_t)=\ln(1+R)-\ln(1+B)=k(R-B)$.
\end{proof}

\begin{table}[!tb]
\centering
\caption{Attribution of an illustrative industry portfolio against the cap-weighted ten industries,
January to \atLast, in percentage points of cumulative return, linked over months by
\cref{prop7:carino}. Weights are those of January, in percent. The benchmark weights drift from month
to month, so industries held at their January weight still have allocation and interaction effects. Entries are
rounded, so the weights shown add to \atWSum{} and the three effects to \atEffSum, against an active
return of \atActive.}
\label{tab7:attrib}
\small
\begin{tabular}{lrrrrr}
\toprule
 & benchmark & portfolio & allocation & selection & interaction\\
 & weight & weight & & & \\
\midrule
\atTableRows
\bottomrule
\end{tabular}
\end{table}

These are the factors of \textcite{carino1999}. As an example, take the cap-weighted ten industries
as the benchmark from January to \atLast, with the market values in the file for each month, and an
illustrative portfolio that starts from the benchmark weights of January, holds ten points less in
technology and three or four points more in health care, energy and utilities, and is rebalanced to
those weights every month. Because the benchmark weights drift with prices, and in July change when the
library reassigns firms to industries, while the portfolio's are reset, industries held at their January
weight still show allocation and interaction effects, and large ones for industries such as
manufacturing whose returns moved far from the benchmark's in single months. Within each industry the portfolio holds the equal-weighted industry
return in place of the value-weighted one, so its selection effect is a bet on smaller companies.
\Cref{tab7:attrib} gives the result. The portfolio returned \atP\% against \atB\% for the benchmark,
an active return of \atActive{} points; the \atMonthsWord{} monthly active returns add to only
\atSumMonthly. After linking, allocation contributed \atAlloc{} points, mostly from the overweight in
energy, selection \atSel{} and interaction \atInter. The linking factors $k_t/k$ ran from \atKmin{} to
\atKmax.
The two functions below implement \cref{eq7:brinson} and \cref{prop7:carino}. Run as a script they
reproduce the answers to \cref{ex7:six} and \cref{ex7:seven}, and \texttt{ch07lib} uses the same
formulas for \cref{tab7:attrib}.
\Needspace*{28\baselineskip}
\lstinputlisting[language=Python]{code/ch07/snippets/brinson_carino.py}

\FloatBarrier
% =====================================================================================
\section{Black-Litterman: starting from the market}
\label{sec7:bl}

\Cref{sec7:estimation} showed that sample means make poor inputs. \textcite{black1992} start instead
from the means that would make the market portfolio optimal. If an investor with risk aversion
$\delta$ chooses $w=(\delta\Sigma)^{-1}\mu_x$, the excess means that make the market weights $w_m$
optimal are the equilibrium returns
\begin{equation}\label{eq7:pi}
  \pi=\delta\Sigma w_m .
\end{equation}
Views are statements $P\mu_x=q+e$ about portfolios $P$ of the assets, with errors $e\sim\N(0,\Omega)$,
and the prior is $\mu_x\sim\N(\pi,\tau\Sigma)$.

\begin{figure}[!tb]
  \centering
  \includegraphics[width=0.94\textwidth]{ch07/fig10_bl}
  \caption{Weights of the ten industries at the end of \fqLast. Left, market weights and the
  Black-Litterman weights with the two views of \cref{sec7:bl}; only the three industries named in
  the views move. Right, unconstrained mean-variance weights from the sample means of the last
  \blW{} months; note the different scale.}
  \label{fig7:bl}
\end{figure}

\begin{proposition}[The Black-Litterman posterior]\label{prop7:bl}
The posterior mean of $\mu_x$ given the views is
\begin{equation}\label{eq7:bl}
  \mu_{\mathrm{BL}}=\bigl[(\tau\Sigma)^{-1}+P^\top\Omega^{-1}P\bigr]^{-1}\bigl[(\tau\Sigma)^{-1}\pi+P^\top\Omega^{-1}q\bigr]
  =\pi+\tau\Sigma P^\top\bigl(P\tau\Sigma P^\top+\Omega\bigr)^{-1}(q-P\pi),
\end{equation}
and the weights it implies are
\begin{equation}\label{eq7:blw}
  w_{\mathrm{BL}}=(\delta\Sigma)^{-1}\mu_{\mathrm{BL}}=w_m+P^\top\Lambda,\qquad
  \Lambda=\frac\tau\delta\bigl(P\tau\Sigma P^\top+\Omega\bigr)^{-1}(q-P\pi).
\end{equation}
\end{proposition}

\begin{proof}
The prior density is proportional to $\exp\{-\frac12(\mu-\pi)^\top(\tau\Sigma)^{-1}(\mu-\pi)\}$ and the
likelihood of the views to $\exp\{-\frac12(P\mu-q)^\top\Omega^{-1}(P\mu-q)\}$. Their product is Gaussian in
$\mu$ with precision $(\tau\Sigma)^{-1}+P^\top\Omega^{-1}P$, and completing the square gives the first
form. For the second, write $G=\tau\Sigma$ and $M=PGP^\top+\Omega$, and use
$(G^{-1}+P^\top\Omega^{-1}P)^{-1}=G-GP^\top M^{-1}PG$, which \cref{prop7:woodbury} gives with the roles
of the blocks exchanged. Multiplying by $G^{-1}\pi+P^\top\Omega^{-1}q$ and using $PGP^\top=M-\Omega$ leaves
$\pi+GP^\top M^{-1}(q-P\pi)$. Finally $(\delta\Sigma)^{-1}\pi=w_m$ by \cref{eq7:pi}, and
$(\delta\Sigma)^{-1}\tau\Sigma P^\top=(\tau/\delta)P^\top$.
\end{proof}

The second statement, due to \textcite{he2002}, is the reason the model behaves well: the investor
departs from the market only along the portfolios that carry views, by amounts that grow with the
surprise $q-P\pi$ and shrink with the uncertainty $\Omega$. Assets on which there is no view keep their
market weights, whatever the estimated covariances. The weights are unconstrained mean-variance demands
and need not add to one: a view on an absolute return changes the total risky exposure, and the
balance is borrowed or lent at the riskless rate. The listing below computes \cref{eq7:bl} and
\cref{eq7:blw} for the two assets of \cref{ex7:eight}.
\Needspace*{15\baselineskip}
\lstinputlisting[language=Python]{code/ch07/snippets/black_litterman.py}

For the ten industries at the end of \fqLast, take $\Sigma$ from the last \blW{} months, the market
weights of the file, $\tau=\blTauFrac$, and $\delta=\blDelta$, the ratio of the long-run mean of the
market excess return to its variance. The equilibrium returns run from \blPiMin\% a year for
\blPiMinName{} to \blPiMax\% for \blPiMaxName, with \blPiHiTec\% for technology, whose weight is
\blWmHiTec\%. Two views: health care will beat technology by 2\% a year, and energy will earn 5\% over
bills. With $\Omega$ equal to the diagonal of $P\tau\Sigma P^\top$, each view is held with the same
confidence as the prior. The gap between health care and technology moves from \blViewOneBefore\% under
the prior to \blViewOneAfter\% and energy from \blViewTwoBefore\% to \blViewTwoAfter\%. The weights move
by $\Lambda=(\blLamOne,\blLamTwo)$ along the two view portfolios: health care from \blWmHlth\% to
\blWblHlth\%, technology from \blWmHiTec\% to \blWblHiTec\%, energy from \blWmEnrgy\% to \blWblEnrgy\%,
and every other industry stays at its market weight (\cref{fig7:bl}). The first view portfolio is long
and short in equal amounts, while the second adds \blLamTwo{} of capital to energy, so the
Black-Litterman weights add to \blBlSum{}: the excess over one is borrowed at the riskless rate, and the
weights are not rescaled to be fully invested. Mean-variance weights from the
\blW-month sample means, with the same $\delta$ and $\Sigma$, have gross exposure of \blMvGross{} times
capital and a short position of \blMvMax\% in \blMvMaxName.



\FloatBarrier
% =====================================================================================
\section{Risk budgets and rebalancing}
\label{sec7:risk}

Institutional allocators often set targets for the share of risk each asset class contributes
instead of its share of capital. Portfolio volatility $\sigma_p(w)=\sqrt{w^\top\Sigma w}$ is
homogeneous of degree one in $w$, so by Euler's theorem
\begin{equation}\label{eq7:euler}
  \sigma_p=\sum_iw_i\frac{\partial\sigma_p}{\partial w_i}=\sum_i\mathrm{RC}_i,\qquad
  \mathrm{RC}_i=\frac{w_i(\Sigma w)_i}{\sigma_p}.
\end{equation}
The risk contribution $\mathrm{RC}_i$ is the part of the volatility attributable to asset $i$, in units
of volatility, and $\mathrm{RC}_i/\sigma_p$ is its share; both can be negative for a hedge.

\begin{proposition}[Equal risk contributions]\label{prop7:erc}
There is exactly one fully invested portfolio with positive weights whose risk contributions are
equal. It is $y/\mathbf 1^\top y$, where $y$ is the unique minimiser over $y>0$ of
$\frac12y^\top\Sigma y-\sum_i\ln y_i$.
\end{proposition}

\begin{proof}
The objective is strictly convex on the open positive orthant and tends to infinity at its boundary
and at infinity, so it has a unique minimiser, where $(\Sigma y)_i=1/y_i$, that is $y_i(\Sigma y)_i=1$
for every $i$. Multiplying the weights by a positive number multiplies every product $w_i(\Sigma w)_i$ by its square
and every contribution $\mathrm{RC}_i$, which divides those products by $\sigma_p$, by the number itself,
so $y/\mathbf 1^\top y$ also has equal contributions. Conversely, if $w>0$ has equal contributions, then
$w_i(\Sigma w)_i=c$ for all $i$, and then $y=w/\sqrt c$
satisfies the first-order condition, so $y$ is the minimiser and $w$ is its normalisation.
\end{proof}

\Cref{fig7:risk} applies this to the market and the five long-short factor premia of the data
library, treating each as an asset class with its own risk premium, as the alternative risk premia
funds do. With equal notional weights the market supplies \rpEqMktShare\% of the risk with one sixth of the
capital. The equal-contribution
weights run from \rpErcMin\% in \rpErcMinName{} to \rpErcMax\% in \rpErcMaxName, with \rpErcMkt\% in the
market. Over the whole sample the equal-contribution portfolio has a Sharpe ratio of \rpSRerc, against \rpSRmkt{} for the market, and with
weights estimated from the previous \csW{} months it is \rpSRrolling. The volatility is only \rpVolErc\% a
year, so matching the market's volatility needs leverage of \rpLev{} times. These are returns of paper
portfolios before trading costs, built by researchers who knew the history, and their published Sharpe
ratios are likely to overstate what a fund could earn after costs, a point \textcite{harvey2016} make forcefully.
\textcite{maillard2010} prove \cref{prop7:erc} and study how such portfolios compare with minimum
variance and equal weights.
The Octave listing solves the first-order condition $y_i(\Sigma y)_i=1$ of the proof with
\texttt{fsolve}, checks the answer to \cref{ex7:nine}, and reproduces these weights from the sample
covariance of the six premia.
\Needspace*{16\baselineskip}
\lstinputlisting[language=Octave]{code/ch07/snippets/erc.m}

\begin{figure}[!tb]
  \centering
  \includegraphics[width=0.84\textwidth]{ch07/fig11_risk}
  \caption{Shares of portfolio volatility, \cref{eq7:euler}, contributed by the market and the five
  factor premia, \fqFirst{} to \fqLast, with equal weights and with equal risk contributions. Labels
  give the equal-contribution weights.}
  \label{fig7:risk}
\end{figure}


\paragraph{Rebalancing.}
A portfolio rebalanced to fixed weights sells what has risen and buys what has fallen, and, to
second order, its expected compound growth rate exceeds the weighted average of the expected growth
rates of its assets.

\begin{proposition}[Diversification return]\label{prop7:rebal}
Let a portfolio be rebalanced each period to weights $w\ge0$ that add to one, and let asset $i$ have
simple return $R_i$ with mean $\mu_i$ and variance $\sigma_i^2$. To second order in the returns,
\begin{equation}\label{eq7:rebal}
  \E\ln(1+R_p)-\sum_iw_i\E\ln(1+R_i)\approx\tfrac12\Bigl(\sum_iw_i\sigma_i^2-\sigma_p^2\Bigr)
  +\tfrac12\Bigl(\sum_iw_i\mu_i^2-\mu_p^2\Bigr),
\end{equation}
and both brackets are nonnegative.
\end{proposition}

\begin{proof}
To second order $\E\ln(1+R)\approx\mu-\frac12(\sigma^2+\mu^2)$, and $\mu_p=\sum_iw_i\mu_i$. The second
bracket is nonnegative by convexity of the square. For the first, correlations are at most one, so
$\sigma_p^2\le\bigl(\sum_iw_i\sigma_i\bigr)^2$, and $\bigl(\sum_iw_i\sigma_i\bigr)^2\le\sum_iw_i\sigma_i^2$
by convexity again.
\end{proof}

For the ten industries rebalanced monthly to equal weights from \fqFirst, the growth rate of
total returns is \rbGrowReb\% a year, against \rbGrowAvg\% for the average of the industries' own growth rates. The
difference, \rbDivActual{} points before rounding, is matched by the first bracket of \cref{eq7:rebal}, \rbDivApprox;
the second is negligible. A portfolio bought in equal weights and never rebalanced grew at
\rbGrowHold\%, between the two, and ended with \rbHoldTop\% in \rbHoldTopName. The gain from rebalancing is
real but it is not free: it comes from the variance of the assets, and it can be offset by
trends that a buy-and-hold portfolio rides. \textcite{booth1992} give the decomposition.

\FloatBarrier
% =====================================================================================
\section{Case study: would estimating the frontier have beaten equal weights?}
\label{sec7:case}

The previous sections give reasons to doubt that an estimated frontier is worth much out of sample,
and one reason to hope: covariances are estimated far better than means. The case study puts this to a
test in the spirit of \textcite{demiguel2009}. At the start of each month from \csStart{} to \csEnd,
\csTT{} months in all, an investor uses only the previous \csW{} months of excess returns of the ten
industries to form seven portfolios, holds them for the month and records the excess return. The
rules are equal weights; market weights; the minimum-variance portfolio of \cref{prop7:frontier}
from the sample covariance; the same from a covariance shrunk towards a multiple of the identity by
the rule of \textcite{ledoit2004}; the long-only minimum-variance portfolio; the direction of the
sample tangency portfolio, scaled to the in-window volatility of equal weights with the rest in bills;
and the long-only tangency portfolio of \cref{prop7:longonly}. In \csFallback{} months every window
mean was negative and the long-only tangency rule held the long-only minimum-variance portfolio
instead. Sharpe ratios are compared with that of equal weights by the test of \textcite{jobson1981}
with the correction of \textcite{memmel2003}, whose statistic is approximately standard normal when
the two Sharpe ratios are equal. Turnover is the average monthly sum of absolute changes in weights
after the drift of the previous month. The returns are before transaction costs and financing
frictions, and the scaled tangency rule can need large long and short positions.

\begin{table}[t]
\centering
\caption{Out-of-sample excess returns of seven rules for the ten industries, \csStart{} to \csEnd,
with weights estimated from the previous \csW{} months. Means, volatilities and alphas are in percent
a year, Sharpe ratios annualised, $z$ is the Memmel statistic for the difference from the Sharpe ratio
of equal weights, turnover is in percent a month, and $\alpha$ and $\beta$ are from a regression on the
market factor.}
\label{tab7:oos}
\small
\begin{tabular}{lrrrrrrr}
\toprule
 & mean & volatility & Sharpe & $z$ & turnover & $\alpha$ & $\beta$\\
\midrule
\csTableRows
\bottomrule
\end{tabular}
\end{table}

\Cref{tab7:oos} and \cref{fig7:case} give the results. Equal weights earned a Sharpe ratio of
\csEwSR. Market weights did slightly worse at \csCapSR, and the three minimum-variance rules matched or
slightly beat equal weights, at \csGmvSR, \csLwSR{} and \csLogmvSR, with the shrunk covariance the best of the
seven. None of
these differences is significant: the $z$ statistics run from \csCapZ{} for market weights to
\csLwZ{} for the shrunk rule. The $z$ statistics
are approximate, since the Memmel test treats monthly returns as independent draws from a
normal distribution, while the weights of each rule carry over from month to month. The rules that use
means did worse. The long-only tangency portfolio matched market weights at \csLotanSR, and the scaled
sample tangency portfolio earned only \csTanSR, with a monthly turnover of \csTanTurn\% and a $z$
statistic of \csTanZ. Estimated means cost the investor most of the Sharpe ratio, as
\cref{prop7:bias} warned, and estimated covariances neither helped nor hurt much.

\begin{figure}[!tb]
  \centering
  \includegraphics[width=0.94\textwidth]{ch07/fig12_case}
  \caption{Top, cumulative log excess returns of five of the rules of \cref{tab7:oos}, \csStart{} to
  \csEnd. Bottom, the Sharpe ratio of four rules over trailing ten-year windows, minus the Sharpe ratio of
  equal weights over the same windows.}
  \label{fig7:case}
\end{figure}

The minimum-variance rules did earn something that the Sharpe ratios hide. Their betas on the market
were between \csGmvB{} and \csLogmvB{}, and their CAPM alphas, \csGmvA\%, \csLwA\% and \csLogmvA\% a
year, have t statistics of \csGmvAt, \csLwAt{} and \csLogmvAt. That is the flat security market line of
\cref{fig7:sml} again: low-beta portfolios earned more than the CAPM says. Levered to a beta of one by
borrowing at the bill rate, the three rules would have had mean excess returns of \csGmvLev\%,
\csLwLev\% and \csLogmvLev\% against \csMktMean\% for the market factor, before any cost of borrowing
above that rate and at a much higher volatility. Leverage scales volatility with the mean, though, so the Sharpe ratios of
\cref{tab7:oos} do not change and at equal volatility the gain over equal weights stays small; much of
these portfolios' variance, \csGmvRes\%, \csLwRes\% and \csLogmvRes\% for the three rules, is residual
risk that a low beta does not remove. With windows of 120 months the ranking is the same: equal weights \csTenEwSR, the shrunk
minimum-variance portfolio \csTenLwSR{} with $z=\csTenLwZ$, the long-only tangency \csTenLotanSR, and the
scaled tangency \csTenTanSR{} with $z=\csTenTanZ$. The median shrinkage intensity was \csShrinkMed, so
even a little shrinkage of the covariance is enough to stabilise the weights.

The core of the experiment is short. The listing below runs the two minimum-variance rules with the
shrinkage of \textcite{ledoit2004} written out in NumPy, and prints the Sharpe ratios \csGmvSR{} and
\csLwSR{} of \cref{tab7:oos}; other universes need only a different data file.
\Needspace*{12\baselineskip}
\lstinputlisting[language=Python]{code/ch07/snippets/rolling_gmv.py}


The answer to the question in the title is no. Over more than half a century, the frontier
estimated from five years of data did not beat equal weights by a margin the Sharpe ratio test can detect, and the
version that trusted estimated means lost badly. For an allocator the implication is about where to
spend effort. A covariance matrix, shrunk and used for minimum variance or for risk budgets, is a
reasonable input. Expected returns should come from somewhere other than the recent past: from
equilibrium, as in \cref{sec7:bl}, or from views whose uncertainty is stated.

\FloatBarrier
% =====================================================================================
\section{Limits of the models}
\label{sec7:limits}

Mean-variance analysis ranks portfolios by two moments. Returns of alternative investments are often
skewed and fat tailed, and an investor who cares about drawdowns or about the left tail will not be
indifferent between two portfolios with the same mean and variance; Chapter~8 treats measures of tail
risk. The analysis is also for one period. A pension plan with liabilities cares about the surplus
of assets over liabilities, and the frontier of surplus returns differs from the frontier of asset
returns because the liability has its own duration, the subject of the immunisation in Chapter~6.

The data here are liquid portfolios of listed stocks with monthly returns. For private funds, the
returns are smoothed, the histories short and the samples selected, so every estimation problem of
\cref{sec7:estimation,sec7:alpha} is worse and \cref{prop7:stale} has to be applied before anything
else. The factor premia in \cref{sec7:risk} are paper portfolios before costs. The CAPM and the
arbitrage pricing theory rest on assumptions, common beliefs, frictionless trading and no arbitrage,
that are idealisations. The value of the theory here is as a set of identities. \cref{prop7:capm} and
\cref{prop7:grs} are algebra: they hold for any means and nonsingular covariance matrix that meet the
conditions stated in them, population or sample, and say exactly what a test of beta pricing is
testing; reading them as equilibrium statements, and the exact law of the joint test,
need further assumptions. Least squares
alphas assume constant betas, which \cref{fig7:rolling} contradicts, and linear payoffs, which
\cref{sec7:options} contradicts. The case study uses one universe and one family of rules, and other
universes, especially those with many assets and few months, favour shrinkage more strongly.

\FloatBarrier
% =====================================================================================
\section{Summary}
\label{sec7:summary}

Equal weights spread over more of a fixed set of assets remove variance, on average over the choice
of assets, down to the average covariance and no further; optimised weights can go lower. The minimum-variance
frontier has a closed form in four scalars, every frontier portfolio is a mix of two, and with a
riskless asset at which all can borrow and lend, every mean-variance investor with the same means
and covariances holds the tangency portfolio in some proportion, whose squared Sharpe ratio is
$\mu_x^\top\Sigma^{-1}\mu_x$. Long-only portfolios solve a convex problem whose solution is a frontier
portfolio of the assets held. Estimated from data, the maximum Sharpe ratio is biased up by a known amount and
the estimated portfolio's true Sharpe ratio is much lower. Beta pricing against a portfolio is
equivalent to that portfolio being efficient, so the CAPM is a claim about the market portfolio. Factor
models make covariances estimable and, with no arbitrage, make the intercept on traded factors the
measure of alpha. The standard error of alpha is that of a mean of residuals, so a good information
ratio takes well over a decade to detect, and screening many managers needs a false discovery rule.
The joint test of alphas is a test of whether the assets improve the factors' Sharpe ratio; for the
ten industries it rejects zero alphas for the six-factor model and not for the CAPM. In the smoothing model, stale prices hide
beta and manufacture alpha, and summing slopes over enough lags recovers the beta. Option-like strategies need a benchmark of options.
Attribution splits active return exactly, and linking factors make it add over time. Black-Litterman
moves away from market weights only along the views. Equal risk contributions have a unique solution,
and, to second order, rebalancing earns about half the gap between average variance and portfolio variance. Out of sample, no
estimated frontier beat equal weights by a margin a test can detect, and estimated means did harm.

\FloatBarrier
% =====================================================================================
\section{Exercises}
\label{sec7:exercises}

\begin{enumerate}[label=\textbf{7.\arabic*}, leftmargin=*, ref=7.\arabic*]
  \item\label{ex7:one} Two assets have volatilities of 20\% and 10\% and a correlation of 0.3. Find
  the weight in the first asset of the minimum-variance portfolio and its volatility.
  \item\label{ex7:two} Three uncorrelated assets have excess means of 4\%, 6\% and 8\% and
  volatilities of 10\%, 15\% and 20\%. Find the fully invested tangency portfolio, its mean,
  volatility and Sharpe ratio, and check the Sharpe ratio against \cref{eq7:theta}.
  \item\label{ex7:three} Compute the beta of each asset of \cref{ex7:two} on its tangency portfolio and
  check \cref{prop7:capm}.
  \item\label{ex7:four} A manager has an annual information ratio of 0.4 against a market with an
  annual Sharpe ratio of 0.45. Using monthly data, how many years does it take for the expected t
  statistic of alpha to reach two?
  \item\label{ex7:five} A private fund's true returns have a beta of 1.2, and it reports returns
  smoothed with $\phi=0.5$. The market premium is 6\% a year. What beta and what spurious alpha does a
  regression on the current market return find, and what does the sum of slopes on the current and two
  lagged market returns give?
  \item\label{ex7:six} A benchmark holds 60\% in equities returning 10\% and 40\% in bonds returning 2\%.
  A portfolio holds 70\% in equities returning 12\% and 30\% in bonds returning 1\%. Split the active
  return into allocation, selection and interaction.
  \item\label{ex7:seven} A portfolio returns 2\% and then 3\%, and its benchmark 1\% in each month.
  Compute the cumulative active return, the sum of the monthly active returns and the linked monthly
  effects of \cref{prop7:carino}.
  \item\label{ex7:eight} Two uncorrelated assets have volatilities of 20\% and 30\% and market weights
  of one half each. With $\delta=2.5$ and $\tau=0.05$, find the equilibrium returns, then the posterior
  mean and the weights after a view that the first asset will earn 3\% with $\Omega=\tau\sigma_1^2$.
  \item\label{ex7:nine} Two assets have volatilities of 20\% and 5\% and a correlation of 0.2. Show that
  the equal-risk-contribution weights do not depend on the correlation and find them.
\end{enumerate}

\subsection*{Answers}

\begin{enumerate}[label=\textbf{7.\arabic*}, leftmargin=*]
  \item The weight is $(\sigma_2^2-\rho\sigma_1\sigma_2)/(\sigma_1^2+\sigma_2^2-2\rho\sigma_1\sigma_2)
  =0.004/0.038=\xrOneW\%$, and the volatility is \xrOneVol\%.
  \item With a diagonal covariance the weights are proportional to $\mu_i/\sigma_i^2$, that is to 4,
  2.67 and 2, so they are \xrTwoA\%, \xrTwoB\% and \xrTwoC\%. The mean is \xrTwoMu\%, the volatility
  \xrTwoVol\% and the Sharpe ratio \xrTwoSR, which equals $\sqrt{0.4^2+0.4^2+0.4^2}$.
  \item The betas are \xrThreeA, \xrThreeB{} and \xrThreeC, and multiplying each by the tangency mean of
  \xrTwoMu\% returns 4\%, 6\% and 8\%.
  \item By \cref{eq7:detect}, $T^*=4(1+0.45^2/12)/(0.4^2/12)$ months, which is \xrFour{} years.
  \item The slope is $(1-\phi)\beta=\xrFiveB$, the spurious alpha $0.5\times1.2\times6\%=\xrFiveA\%$ a
  year, and the sum of three slopes $\beta(1-\phi^3)=\xrFiveD$.
  \item The benchmark returns \xrSixRb\% and the portfolio $8.70\%$, an active return of \xrSixActive{}
  points. Allocation is $0.1(10-6.8)-0.1(2-6.8)=\xrSixAlloc$, selection $0.6\times2-0.4\times1=\xrSixSel$
  and interaction $0.1\times2+0.1\times1=\xrSixInt$.
  \item The portfolio compounds to \xrSevenP\% and the benchmark to \xrSevenB\%, an active return of
  3.05 points, while the monthly active returns add to \xrSevenSum. The scaled factors are
  \xrSevenKone{} and \xrSevenKtwo, so the linked effects are \xrSevenOne{} and \xrSevenTwo, which add
  to 3.05.
  \item The equilibrium returns are $\pi=2.5\Sigma w_m$, \xrEightPiA\% and \xrEightPiB\%. With
  $\Omega=\tau\sigma_1^2$ the view and the prior are equally precise, so the posterior mean of the first
  asset is halfway, \xrEightMuA\%, and the second is unchanged. The weights are \xrEightWA\% and
  \xrEightWB\%, a move along the view portfolio only. The weights add to 90\%; the rest is in bills, since
$(\delta\Sigma)^{-1}\mu_{\mathrm{BL}}$ need not be fully invested.
  \item With two assets the contributions are equal when $w_1^2\sigma_1^2+w_1w_2\rho\sigma_1\sigma_2=w_2^2\sigma_2^2+w_1w_2\rho\sigma_1\sigma_2$,
  so $w_1\sigma_1=w_2\sigma_2$ whatever $\rho$, and $w_1=\sigma_2/(\sigma_1+\sigma_2)=\xrNineW\%$. The
  portfolio volatility is \xrNineVol\%.
\end{enumerate}

\FloatBarrier
% =====================================================================================
\section{Notes, sources and credits}
\label{sec7:credits}

\paragraph{Sources}
Portfolio selection by mean and variance is from \textcite{markowitz1952}, the closed form of the
frontier from \textcite{merton1972}, and the separation of the riskless asset from the risky portfolio
from \textcite{tobin1958}. The CAPM is from \textcite{sharpe1964} and \textcite{lintner1965}, the
zero-beta version from \textcite{black1972}, and the critique of its tests from \textcite{roll1977}. The
arbitrage pricing theory is from \textcite{ross1976} and its approximate form from \textcite{huberman1982}; the factors used here are from
\textcite{fama1993,fama2015} and \textcite{carhart1997}. The bias of estimated portfolios is studied by
\textcite{michaud1989} and \textcite{kan2007}, and the Wishart moments are in \textcite{muirhead1982}.
The joint test of alphas is from \textcite{gibbons1989}, the standard errors from \textcite{newey1987},
the false discovery rule from \textcite{benjamini1995} and \textcite{storey2002}, and its application to
funds and factors from \textcite{barras2010} and \textcite{harvey2016}. Betas of thinly traded assets are
from \textcite{dimson1979} and smoothing in hedge fund returns from \textcite{getmansky2004}. Market
timing regressions are from \textcite{treynor1966} and \textcite{henriksson1981}, the option-based
benchmark from \textcite{glosten1994}, and option-like hedge fund exposures from \textcite{agarwal2004}.
Attribution is from \textcite{brinson1985,brinson1986} and linking from \textcite{carino1999}. The
Black-Litterman model is from \textcite{black1992} and its interpretation from \textcite{he2002}. Equal
risk contributions are from \textcite{maillard2010}, diversification returns from \textcite{booth1992},
shrinkage from \textcite{ledoit2004}, the out-of-sample comparison from \textcite{demiguel2009}, and the
Sharpe ratio test from \textcite{jobson1981} and \textcite{memmel2003}.

\paragraph{Software used}
The scripts call NumPy \parencite{numpy} for arrays, linear algebra and random numbers, SciPy \parencite{scipy} for
distributions, pandas \parencite{pandas} to read the data files, and Matplotlib \parencite{matplotlib} for the figures; the
tests compare the regressions with statsmodels \parencite{statsmodels}, which the listing of
\cref{sec7:alpha} also uses. The cross-checks run in R \parencite{rcore} and
GNU Octave \parencite{octave}. \texttt{code/requirements-lock.txt} records the full environment.

\paragraph{Projects for comparison}
The code here does not call the following projects, but they implement the same methods and are
useful for checking: PyPortfolioOpt \parencite{pyportfolioopt} has mean-variance optimisation, shrinkage estimators and
Black-Litterman; Riskfolio-Lib \parencite{riskfolio} has risk budgeting and many other risk measures;
\texttt{linearmodels} \parencite{linearmodels} has factor model estimation with the joint test of alphas; and
\texttt{pandas-datareader} \parencite{datareader} downloads the files of the data library used here. The full list, with links, is
in \texttt{CREDITS.md}.

\paragraph{Data}
The factors and industry portfolios are from the data library of \textcite{frenchlib}, in the release
built from the CRSP database of August 2026 and retrieved on 19 October 2026. The files are not redistributed with the code:
\texttt{scripts/build\_french.py -{}-download} saves the CSV files from the library's zip archives
in \texttt{data/raw/french/} and builds the two files the code reads. The library is revised from time
to time, and the numbers here use the saved files. The put-writing fund of \cref{sec7:options}, the
managers of \cref{sec7:alpha} and the portfolio of \cref{tab7:attrib} are illustrations and describe no
real fund.

\printbibliography[heading=subbibliography,title={References for Chapter \thechapter}]
\end{refsection}
